% Support & locomotion — Biology Upper Sixth — Cours signé OnBuch+
% Official MINESEC syllabus. Compiler avec : tectonic BIO05-support-locomotion.tex
\newcommand{\DOCMATIERE}{Biology}
\newcommand{\DOCNIVEAU}{Upper Sixth}
\newcommand{\DOCMODULE}{Upper Sixth — Biology}
\newcommand{\DOCLECON}{5}
\newcommand{\DOCTITRE}{Support \& locomotion}
% OnBuch+ — préambule « cours signés » ENRICHI (compilé avec tectonic / XeLaTeX)
% Système visuel « Pop » d'OnBuch+ : fond crème (jamais blanc), bordures encre
% épaisses, ombres « dures » décalées, titres Archivo Black en capitales.
% Version MATHÉMATIQUES : graphes (pgfplots/tikz), boîtes pédagogiques
% (définition, propriété, méthode, attention, le savais-tu, exercice résolu,
% exercice, TP, bilan), page de garde et signature OnBuch+ ancrée en bas.
%
% Macros attendues dans main.tex AVANT \input{preamble} :
%   \DOCMATIERE  \DOCNIVEAU  \DOCMODULE  \DOCLECON (numéro)  \DOCTITRE
\documentclass[11pt]{article}

\usepackage{fontspec}
\usepackage[english]{babel}
\usepackage[a4paper,margin=2cm,top=3.4cm,bottom=2.8cm,headheight=1.4cm,footskip=1.3cm]{geometry}
\usepackage{amsmath,amssymb,amsfonts}
\usepackage{mathtools} % \xmapsto, \coloneqq... souvent utilisés par les modèles
\usepackage{cancel} % simplifications barrées (bilans de forces)
% \abs{z} (module, valeur absolue) et \norm{u} (norme), utilisés par les modèles
\providecommand{\abs}{}\let\abs\relax\DeclarePairedDelimiter{\abs}{\lvert}{\rvert}
\providecommand{\norm}{}\let\norm\relax\DeclarePairedDelimiter{\norm}{\lVert}{\rVert}
\usepackage[table]{xcolor} % [table] : fournit aussi \rowcolors (alternance de lignes)
\usepackage{pagecolor}
\usepackage{tikz}
\usetikzlibrary{arrows.meta,positioning,calc,shapes.geometric,shapes.misc,shapes.arrows,decorations.pathmorphing,decorations.pathreplacing,decorations.markings,patterns,shadows,mindmap,trees,fit,backgrounds,matrix,intersections,angles,quotes}
\usepackage{pgfplots}
\usepackage[version=4]{mhchem} % équations nucléaires \ce{^{235}_{92}U}
\usepackage[european,siunitx]{circuitikz} % schémas électriques
\pgfplotsset{compat=1.18}
\usepgfplotslibrary{fillbetween,groupplots} % aires entre courbes (calcul intégral)
\usepackage{enumitem}
\usepackage{fancyhdr}
% [most] charge toutes les bibliothèques tcolorbox (listings, minted, raster,
% documentation...) et gonfle inutilement la mémoire TeX sur un document long
% (~300 boîtes) : on ne charge que ce qui est réellement utilisé.
\usepackage[skins,breakable]{tcolorbox}
\usepackage{graphicx}
\usepackage{colortbl}
\usepackage{booktabs}
\usepackage{tabularx}
\usepackage{diagbox} % case d'en-tête diagonale des tableaux à double entrée
% Colonnes extensibles centrée / gauche / droite (C, L, R), souvent utilisées par les modèles
\newcolumntype{C}{>{\centering\arraybackslash}X}
\newcolumntype{L}{>{\raggedright\arraybackslash}X}
\newcolumntype{R}{>{\raggedleft\arraybackslash}X}
\usepackage{array}
\usepackage{multicol}
\usepackage{siunitx}
% parse-numbers=false : accepte aussi \SI{2,5 \times 10^{-3}}{...}
\sisetup{locale=UK,output-decimal-marker={.},group-minimum-digits=4,parse-numbers=false}
\DeclareSIUnit\ohm{\ensuremath{\symup{\Omega}}} % U+2126 absent de la police
\DeclareSIUnit\division{div} % réglages d'oscilloscope (ms/div, V/div)
\DeclareSIUnit\tr{tr} % tours (vitesse de rotation, ex. \SI{1500}{\tr\per\minute})
\usepackage{qrcode}
% \degree : commande fréquemment utilisée par les modèles pour un angle ou une
% température en dehors de \SI{}{} (non fournie par les paquets chargés).
\providecommand{\degree}{^{\circ}}
% \qed (fin de démonstration), utilisé par les modèles sans amsthm
\providecommand{\qed}{\hfill$\square$}
% \barre : double barre des valeurs interdites dans un tableau de variations (tkz-tab)
\providecommand{\barre}{\ensuremath{\|}}
% environnement proof (amsthm non chargé) utilisé par les modèles
\ifdefined\proof\else\newenvironment{proof}[1][Proof]{\par\noindent\textit{#1.}\ }{\qed\par}\fi
\usepackage{lastpage}
\usepackage{hyperref}
\hypersetup{hidelinks}

% ============================================================
% Palette « Pop » OnBuch+ — mêmes hex que la classe P (plus_ui.dart)
% ============================================================
\definecolor{poporange}{HTML}{F59321}
\definecolor{popdark}{HTML}{E07A0C}
\definecolor{popcream}{HTML}{FFF6E8}
\definecolor{popink}{HTML}{1C1714}
\definecolor{poppurple}{HTML}{7A5AE0}
\definecolor{popgreen}{HTML}{2E9E6A}
\definecolor{poppink}{HTML}{E0457B}
\definecolor{popblue}{HTML}{2F7DE1}
\definecolor{popgold}{HTML}{C98A12}
\definecolor{popmuted}{HTML}{8A7F76}
\definecolor{popline}{HTML}{EADFD2}
\definecolor{popgray}{HTML}{8A7F76}
\colorlet{popgrey}{popgray}
% Alias tolérants (le modèle nomme parfois les couleurs différemment)
\colorlet{popwhite}{white}
\colorlet{popblack}{popink}
\colorlet{popred}{poppink}
\colorlet{popyellow}{popgold}
% Teintes claires pour fonds de tableaux / zones
\colorlet{popblueL}{popblue!10!white}
\colorlet{poporangeL}{poporange!14!white}
\colorlet{popgreenL}{popgreen!12!white}
\colorlet{poppurpleL}{poppurple!12!white}
\colorlet{poppinkL}{poppink!12!white}
\colorlet{popgoldL}{popgold!14!white}

\pagecolor{popcream}

% ============================================================
% Typographie
% ============================================================
\defaultfontfeatures{Ligatures=TeX}
\setmainfont{PlusJakartaSans-Medium.ttf}[Path=../../../fonts/,BoldFont=PlusJakartaSans-Bold.ttf]
% Maths en sans-serif (Fira Math) avec chiffres et lettres droites de Plus
% Jakarta, pour que les formules se fondent dans le texte.
\usepackage[math-style=ISO,bold-style=ISO]{unicode-math}
\setmathfont{FiraMath-Regular.otf}[Path=../../../fonts/]
\setmathfont{PlusJakartaSans-Medium.ttf}[Path=../../../fonts/,range={up/num,up/{Latin,latin},\mathup}]
% (icomma retiré : décimales avec point en anglais)
% Caractères mathématiques Unicode absents des polices de texte (Plus Jakarta,
% Archivo Black) : rendus par leur équivalent mathématique, en texte comme en
% formule — sinon ils s'affichent en carré vide (ex. « Arithmétique dans ℤ »).
\usepackage{newunicodechar}
\newunicodechar{ℝ}{\ensuremath{\mathbb{R}}}
\newunicodechar{ℤ}{\ensuremath{\mathbb{Z}}}
\newunicodechar{ℂ}{\ensuremath{\mathbb{C}}}
\newunicodechar{ℕ}{\ensuremath{\mathbb{N}}}
\newunicodechar{ℚ}{\ensuremath{\mathbb{Q}}}
\newunicodechar{𝔻}{\ensuremath{\mathbb{D}}}
\newunicodechar{α}{\ensuremath{\alpha}}
\newunicodechar{β}{\ensuremath{\beta}}
\newunicodechar{γ}{\ensuremath{\gamma}}
\newunicodechar{δ}{\ensuremath{\delta}}
\newunicodechar{ε}{\ensuremath{\varepsilon}}
\newunicodechar{θ}{\ensuremath{\theta}}
\newunicodechar{λ}{\ensuremath{\lambda}}
\newunicodechar{μ}{\ensuremath{\mu}}
\newunicodechar{σ}{\ensuremath{\sigma}}
\newunicodechar{τ}{\ensuremath{\tau}}
\newunicodechar{φ}{\ensuremath{\varphi}}
\newunicodechar{ψ}{\ensuremath{\psi}}
\newunicodechar{ω}{\ensuremath{\omega}}
\newunicodechar{Σ}{\ensuremath{\Sigma}}
% majuscules produites par \MakeUppercase dans les titres (α-aminés -> Α)
\newunicodechar{Α}{\ensuremath{\alpha}}
\newunicodechar{Β}{\ensuremath{\beta}}
\newunicodechar{Γ}{\ensuremath{\Gamma}}
\newunicodechar{Δ}{\ensuremath{\Delta}}
\newunicodechar{Ε}{\ensuremath{\varepsilon}}
\newunicodechar{Θ}{\ensuremath{\Theta}}
\newunicodechar{Λ}{\ensuremath{\Lambda}}
\newunicodechar{Μ}{\ensuremath{\mu}}
\newunicodechar{Τ}{\ensuremath{\tau}}
\newunicodechar{Φ}{\ensuremath{\Phi}}
\newunicodechar{Ψ}{\ensuremath{\Psi}}
\newunicodechar{Ω}{\ensuremath{\Omega}}
\newunicodechar{↦}{\ensuremath{\mapsto}}
\newunicodechar{∈}{\ensuremath{\in}}
\newunicodechar{∉}{\ensuremath{\notin}}
\newunicodechar{∧}{\ensuremath{\wedge}}
\newunicodechar{⇔}{\ensuremath{\Leftrightarrow}}
\newunicodechar{⟺}{\ensuremath{\Longleftrightarrow}}
\newunicodechar{⇒}{\ensuremath{\Rightarrow}}
\newunicodechar{⟹}{\ensuremath{\Longrightarrow}}
\newunicodechar{∘}{\ensuremath{\circ}}
\newunicodechar{∪}{\ensuremath{\cup}}
\newunicodechar{∩}{\ensuremath{\cap}}
\newunicodechar{≡}{\ensuremath{\equiv}}
\newunicodechar{⊂}{\ensuremath{\subset}}
\newunicodechar{⊥}{\ensuremath{\perp}}
\newunicodechar{∃}{\ensuremath{\exists}}
\newunicodechar{∀}{\ensuremath{\forall}}
\newunicodechar{∛}{\ensuremath{\sqrt[3]{\,}}}
\newunicodechar{∜}{\ensuremath{\sqrt[4]{\,}}}
\newunicodechar{⃗}{\ensuremath{{}^{\to}}}
\newunicodechar{ⁿ}{\ifmmode^{n}\else\textsuperscript{\ensuremath{n}}\fi}
\newunicodechar{ˣ}{\ifmmode^{x}\else\textsuperscript{\ensuremath{x}}\fi}
\newunicodechar{ᵇ}{\ifmmode^{b}\else\textsuperscript{\ensuremath{b}}\fi}
\newunicodechar{ʳ}{\ifmmode^{r}\else\textsuperscript{\ensuremath{r}}\fi}
\newunicodechar{ᵘ}{\ifmmode^{u}\else\textsuperscript{\ensuremath{u}}\fi}
\newunicodechar{ʸ}{\ifmmode^{y}\else\textsuperscript{\ensuremath{y}}\fi}
\newunicodechar{ᵃ}{\ifmmode^{a}\else\textsuperscript{\ensuremath{a}}\fi}
\newunicodechar{ᵉ}{\ifmmode^{e}\else\textsuperscript{\ensuremath{e}}\fi}
\newunicodechar{ᵏ}{\ifmmode^{k}\else\textsuperscript{\ensuremath{k}}\fi}
\newunicodechar{ᵖ}{\ifmmode^{p}\else\textsuperscript{\ensuremath{p}}\fi}
\newunicodechar{ᴺ}{\ifmmode^{N}\else\textsuperscript{\ensuremath{N}}\fi}
\newunicodechar{ᶜ}{\ifmmode^{c}\else\textsuperscript{\ensuremath{c}}\fi}
\newunicodechar{ʰ}{\ifmmode^{h}\else\textsuperscript{\ensuremath{h}}\fi}
\newunicodechar{ᵠ}{\ifmmode^{q}\else\textsuperscript{\ensuremath{q}}\fi}
\newunicodechar{ₙ}{\ifmmode_{n}\else\textsubscript{\ensuremath{n}}\fi}
\newunicodechar{ₐ}{\ifmmode_{a}\else\textsubscript{\ensuremath{a}}\fi}
\newunicodechar{ᵢ}{\ifmmode_{i}\else\textsubscript{\ensuremath{i}}\fi}
\newunicodechar{ᵣ}{\ifmmode_{r}\else\textsubscript{\ensuremath{r}}\fi}
\newunicodechar{ₖ}{\ifmmode_{k}\else\textsubscript{\ensuremath{k}}\fi}
\newunicodechar{ᵦ}{\ifmmode_{\beta}\else\textsubscript{\ensuremath{\beta}}\fi}
\newunicodechar{ₓ}{\ifmmode_{x}\else\textsubscript{\ensuremath{x}}\fi}
\newfontfamily\popdisplay{ArchivoBlack-Regular.ttf}[Path=../../../fonts/]
\newfontfamily\poplogo{SpaceGrotesk-Bold.ttf}[Path=../../../fonts/]
\newfontfamily\popsemi{PlusJakartaSans-SemiBold.ttf}[Path=../../../fonts/]
\newfontfamily\popxbold{PlusJakartaSans-ExtraBold.ttf}[Path=../../../fonts/]
\newfontfamily\popxboldls{PlusJakartaSans-ExtraBold.ttf}[Path=../../../fonts/,LetterSpace=3.0]

\setlist[enumerate]{leftmargin=1.7em,itemsep=5pt,topsep=4pt}
\setlist[itemize]{leftmargin=1.7em,itemsep=4pt,topsep=4pt,label={\color{poporange}\textbullet}}
\setlength{\parindent}{0pt}
\setlength{\parskip}{5pt}
\renewcommand{\arraystretch}{1.25}

% pgfplots : style Pop par défaut
\pgfplotsset{
  every axis/.append style={axis line style={popink,line width=1pt},tick style={popink},
    grid style={popline,line width=0.5pt},label style={font=\small},tick label style={font=\footnotesize}},
  popcurve/.style={poporange,line width=1.8pt,no markers,smooth},
}
% Même style disponible dans un \draw TikZ simple (hors pgfplots)
\tikzset{popcurve/.style={poporange,line width=1.8pt}}
\tikzset{popgrid/.style={popline,very thin}}
\colorlet{popcurve}{poporange} % le nom du style est parfois utilisé comme couleur

% ============================================================
% Pill / squiggle
% ============================================================
\newcommand{\poppill}[3]{%
  \tikz[baseline=-0.55ex]\node[draw=popink,line width=1.2pt,rounded corners=2.6pt,%
    fill=#2,inner xsep=6.5pt,inner ysep=3pt]{%
    \popxboldls\fontsize{7}{7}\selectfont\color{#3}\MakeUppercase{#1}};%
}
\newcommand{\popsquiggle}[1]{%
  \tikz[baseline=-0.4ex]{
    \draw[#1,line width=2.6pt,line cap=round]
      plot [smooth] coordinates {(0,0.09) (0.34,0.01) (0.68,0.21) (1.02,0.01) (1.36,0.10)};
  }%
}
% Mot-clé surligné (marqueur orange sous le texte)
\newcommand{\cle}[1]{{\popxbold\color{popdark}#1}}

% ============================================================
% En-tête / pied
% ============================================================
\pagestyle{fancy}
\fancyhf{}
\renewcommand{\headrulewidth}{0pt}
\renewcommand{\footrulewidth}{0pt}
\lhead{\raisebox{-0.15ex}{\poplogo\fontsize{13}{13}\selectfont\color{popink}OnBuch\color{poporange}+}}
\rhead{\poppill{\DOCMATIERE\ \textbullet\ \DOCNIVEAU\ \textbullet\ Chapter \DOCLECON}{white}{popink}}
\lfoot{\popxbold\fontsize{7.6}{7.6}\selectfont\color{popmuted}\MakeUppercase{Signed course OnBuch+}\ \textbullet\ {\popsemi\DOCTITRE}}
\rfoot{\tikz[baseline=-0.6ex]\node[rounded corners=8.5pt,fill=popink,minimum height=17pt,minimum width=17pt,inner xsep=5pt,inner sep=0]{\popxbold\fontsize{8}{8}\selectfont\color{popcream}\thepage\,/\,\pageref*{LastPage}};}

% ============================================================
% Titres
% ============================================================
\newcommand{\chaptitre}[2]{%
  \begin{tcolorbox}[enhanced,colback=poporange,colframe=popink,boxrule=2.6pt,arc=11pt,
      drop shadow=popink,left=15pt,right=15pt,top=12pt,bottom=11pt,before skip=3pt,after skip=18pt]
    \poppill{#2}{popink}{popcream}\par\vspace{9pt}
    {\raggedright\hyphenpenalty=10000\exhyphenpenalty=10000\popdisplay\fontsize{22}{24}\selectfont\color{popcream}\MakeUppercase{#1}\par}
  \end{tcolorbox}
}
% Section numérotée : pastille orange avec le numéro + titre Archivo
\newcounter{popsec}
\newcommand{\coursec}[1]{%
  \stepcounter{popsec}\setcounter{popsub}{0}%
  \par\vspace{16pt}\noindent
  \begin{minipage}{\linewidth}
  \tikz[baseline=-0.7ex]\node[draw=popink,line width=1.6pt,fill=poporange,rounded corners=4pt,minimum size=22pt,inner sep=0]{\popdisplay\fontsize{12}{12}\selectfont\color{popcream}\thepopsec};%
  \hspace{8pt}{\popdisplay\fontsize{15}{17}\selectfont\color{popink}\MakeUppercase{#1}}\par
  \vspace{4pt}\noindent{\color{poporange}\rule{36pt}{3.6pt}}
  \end{minipage}\par\nopagebreak\vspace{8pt}
}
\newcounter{popsub}
\newcommand{\courssub}[1]{%
  \stepcounter{popsub}%
  \par\vspace{9pt}\noindent{\popxbold\fontsize{11.5}{13}\selectfont\color{popink}{\color{poporange}\thepopsec.\thepopsub}\hspace{6pt}#1}\par\nopagebreak\vspace{4pt}
}

% ============================================================
% Boîtes pédagogiques — même recette : fond blanc, bordure épaisse colorée,
% ombre dure encre, badge plein. Titre optionnel [#1].
% ============================================================
\tcbset{popbase/.style={breakable,enhanced,colback=white,boxrule=2.3pt,arc=9pt,
  shadow={3pt}{-3pt}{0pt}{popink},left=14pt,right=14pt,top=11pt,bottom=11pt,before skip=11pt,after skip=15pt,
  fontupper=\popsemi\fontsize{10.6}{14.5}\selectfont\color{popink},
  fontlower=\popsemi\fontsize{10.6}{14.5}\selectfont\color{popink}}}
% Coche dessinée (le glyphe ✓ n'existe pas dans les polices du document)
\newcommand{\popcheck}[1]{\tikz[baseline=-0.5ex]\draw[#1,line width=1.6pt,line cap=round,line join=round](-0.13,0.0)--(-0.04,-0.09)--(0.13,0.1);}
\providecommand{\coche}{\popcheck{popgreen}}
\providecommand{\resultat}[1]{\par\smallskip\noindent\fcolorbox{popgreen}{popgreenL}{\parbox{\dimexpr\linewidth-2\fboxsep-2\fboxrule\relax}{\textbf{Result.} #1}}\par\smallskip}
\newcommand{\popboxhead}[3]{\poppill{#1}{#2}{white}\ifx\relax#3\relax\else\hspace{6pt}{\popxbold\fontsize{10.6}{12}\selectfont\color{popink}#3}\fi\par\vspace{7pt}}

\newtcolorbox{aretenir}[1][]{popbase,colframe=popgreen,before upper={\popboxhead{Key points}{popgreen}{#1}}}
\newtcolorbox{exemplebox}[1][]{popbase,colframe=poporange,before upper={\popboxhead{Exemple}{poporange}{#1}}}
\newtcolorbox{definition}[1][]{popbase,colframe=popblue,before upper={\popboxhead{Definition}{popblue}{#1}}}
\newtcolorbox{propriete}[1][]{popbase,colframe=poppurple,before upper={\popboxhead{Property}{poppurple}{#1}}}
\newtcolorbox{methode}[1][]{popbase,colframe=popgold,before upper={\popboxhead{Method}{popgold}{#1}}}
\newtcolorbox{attention}[1][]{popbase,colframe=poppink,before upper={\popboxhead{Warning}{poppink}{#1}}}
\newtcolorbox{experience}[1][]{popbase,colframe=popblue,colback=popblueL,before upper={\popboxhead{Experiment}{popblue}{#1}}}
% « Le savais-tu ? » : carte violette pleine (contexte, culture, Cameroun)
\newtcolorbox{savaistu}[1][]{popbase,colback=poppurple,colframe=popink,
  fontupper=\popsemi\fontsize{10.4}{14.2}\selectfont\color{white},
  before upper={\poppill{Did you know?}{white}{poppurple}\ifx\relax#1\relax\else\hspace{6pt}{\popxbold\color{white}#1}\fi\par\vspace{7pt}}}
% Exercice résolu : énoncé en haut, \tcblower puis la solution sur fond crème
\newcounter{popexo}\newcounter{popexr}
\newtcolorbox{exoresolu}[1][]{popbase,colframe=poporange,colbacklower=popcream,
  segmentation style={popink,line width=1.2pt,dashed},
  before upper={\stepcounter{popexr}\popboxhead{Worked example \thepopexr}{poporange}{#1}},
  before lower={\poppill{Solution}{popink}{popcream}\par\vspace{6pt}}}
% Exercice (non résolu en place) — la correction est dans la partie Corrigés
\newtcolorbox{exercice}[1][]{popbase,colframe=popink,
  before upper={\stepcounter{popexo}\popboxhead{Exercise \thepopexo}{popink}{#1}}}
% Corrigé
\newtcolorbox{corrige}[1][]{popbase,colframe=popgreen,colback=popgreenL,
  before upper={\popboxhead{Solution}{popgreen}{#1}}}
% Objectifs / prérequis / bilan
\newtcolorbox{objectifs}{popbase,colframe=popink,colback=poporangeL,
  before upper={\popboxhead{Chapter objectives}{popink}{}}}
\newtcolorbox{prerequis}{popbase,colframe=popink,colback=popblueL,
  before upper={\popboxhead{Prerequisites}{popblue}{}}}
\newtcolorbox{bilan}{popbase,colframe=popink,colback=white,boxrule=2.8pt,
  before upper={\popboxhead{Fiche bilan}{poporange}{L'essentiel à connaître pour l'examen}}}
% Figure encadrée avec légende
\newtcolorbox{popfigure}[1][]{enhanced,breakable=false,colback=white,colframe=popline,boxrule=1.4pt,arc=8pt,
  left=10pt,right=10pt,top=10pt,bottom=8pt,before skip=10pt,after skip=12pt,halign=center,
  lower separated=false,halign lower=center,
  fontlower=\popsemi\fontsize{9}{11.5}\selectfont\color{popmuted},
  #1}
\newcommand{\legende}[1]{\par\vspace{6pt}{\popsemi\fontsize{9}{11.5}\selectfont\color{popmuted}#1}}

% ============================================================
% Signature OnBuch+
% ============================================================
\newcommand{\poplogoinline}[1]{{\poplogo\fontsize{#1}{#1}\selectfont\color{popink}OnBuch\color{poporange}+}}
\newcommand{\signatureonbuch}{%
  \begin{tcolorbox}[enhanced,colback=popink,colframe=popink,boxrule=2.2pt,arc=10pt,
      drop shadow=poporange,left=14pt,right=14pt,top=10pt,bottom=10pt,before skip=0pt,after skip=0pt]
    \tikz[baseline=-0.7ex]\node[circle,fill=popgreen,draw=popcream,line width=1.3pt,minimum size=22pt,inner sep=0]{\popcheck{white}};%
    \hspace{9pt}\begin{minipage}[c]{0.62\linewidth}
      {\popxbold\fontsize{12}{13}\selectfont\color{popcream}Signed\ \textbullet\ {\poplogo OnBuch\color{poporange}+}}\par\vspace{2pt}
      {\raggedright\popsemi\fontsize{8}{10.5}\selectfont\color{popline}\MakeUppercase{OnBuch+ teaching team}\\\MakeUppercase{Follows the official MINESEC syllabus}\par}
    \end{minipage}\hfill
    {\poplogo\fontsize{22}{22}\selectfont\color{popcream}OnBuch\color{poporange}+}
  \end{tcolorbox}%
}

% ============================================================
% Page de garde
% #1 titre · #2 matière · #3 niveau · #4 module · #5 n° leçon · #6 durée ·
% #7 description / objectifs (texte ou itemize)
% ============================================================
\newcommand{\pagegarde}[7]{%
  \thispagestyle{empty}
  \newgeometry{a4paper,margin=1.7cm,top=1.6cm,bottom=1.4cm}%
  \begin{tikzpicture}[remember picture,overlay]
    \fill[poporange!18] ([xshift=-3.2cm,yshift=-3.6cm]current page.north east) circle (4.6cm);
    \fill[poppurple!14] ([xshift=2.4cm,yshift=7.6cm]current page.south west) circle (3.8cm);
  \end{tikzpicture}%
  {\poplogo\fontsize{30}{30}\selectfont\color{popink}OnBuch\color{poporange}+}\hfill
  \raisebox{0.6ex}{\poppill{Signed course \textbullet\ Premium}{popink}{popcream}}\par
  \vspace{1pt}\popsquiggle{poppurple}\par
  \vspace{8pt}
  \begin{tcolorbox}[enhanced,colback=poporange,colframe=popink,boxrule=2.8pt,arc=13pt,
      drop shadow=popink,left=18pt,right=18pt,top=12pt,bottom=13pt,after skip=10pt]
    \poppill{#2 \textbullet\ #3}{popink}{popcream}\hspace{5pt}\poppill{#4}{white}{popink}\par\vspace{8pt}
    {\popdisplay\fontsize{14}{14}\selectfont\color{popink}CHAPTER #5}\par\vspace{4pt}
    {\raggedright\hyphenpenalty=10000\exhyphenpenalty=10000\popdisplay\fontsize{25}{27}\selectfont\color{popcream}\MakeUppercase{#1}\par}
    \vspace{8pt}
    {\popxbold\fontsize{9.5}{9.5}\selectfont\color{popink}\MakeUppercase{Suggested time: #6}}
  \end{tcolorbox}
  \begin{tcolorbox}[enhanced,colback=white,colframe=popink,boxrule=2.2pt,arc=10pt,
      drop shadow=popink,left=16pt,right=16pt,top=10pt,bottom=10pt,after skip=8pt]
    \poppill{In this chapter}{poppurple}{white}\par\vspace{5pt}
    \popsemi\fontsize{9.3}{12.6}\selectfont\color{popink}#7
  \end{tcolorbox}
  \noindent\begin{minipage}[t]{0.487\linewidth}
    \begin{tcolorbox}[enhanced,colback=popgreen,colframe=popink,boxrule=2.2pt,arc=10pt,
        drop shadow=popink,left=11pt,right=11pt,top=10pt,bottom=10pt,height=3.1cm,valign=center]
      \tcbox[colback=white,colframe=popink,boxrule=1.3pt,arc=5pt,boxsep=0pt,nobeforeafter,
        left=3pt,right=3pt,top=3pt,bottom=3pt,valign=center]{\qrcode[height=1.9cm]{https://play.google.com/store/apps/details?id=com.luvvix.onbuch&hl=en}}%
      \hspace{8pt}\begin{minipage}[c]{0.5\linewidth}
        {\popdisplay\fontsize{11}{12.5}\selectfont\color{white}\MakeUppercase{Download OnBuch}}\par\vspace{4pt}
        {\raggedright\popsemi\fontsize{8.4}{11}\selectfont\color{white}Free \textbullet\ Google Play\\AI tutor, past papers, results\par}
      \end{minipage}
    \end{tcolorbox}
  \end{minipage}\hfill
  \begin{minipage}[t]{0.487\linewidth}
    \begin{tcolorbox}[enhanced,colback=poppurple,colframe=popink,boxrule=2.2pt,arc=10pt,
        drop shadow=popink,left=11pt,right=11pt,top=10pt,bottom=10pt,height=3.1cm,valign=center]
      \tikz[baseline=-0.7ex]\node[circle,fill=white,draw=popink,line width=1.3pt,minimum size=1.6cm,inner sep=0]{\popdisplay\fontsize{24}{24}\selectfont\color{poppurple}+};%
      \hspace{8pt}\begin{minipage}[c]{0.6\linewidth}
        {\popdisplay\fontsize{11}{12.5}\selectfont\color{white}\MakeUppercase{Go OnBuch+}}\par\vspace{4pt}
        {\raggedright\popsemi\fontsize{8.4}{11}\selectfont\color{white}All signed courses, unlimited AI tutor, PDF sheets — OnBuch+ tab in the app\par}
      \end{minipage}
    \end{tcolorbox}
  \end{minipage}
  \vspace{8pt}
  \popcontenu
  \vfill
  \signatureonbuch
  \restoregeometry
  \newpage
}

% Rangée « Ce cours contient » (page de garde)
\newcommand{\poptile}[3]{% #1 couleur #2 gros texte #3 légende
  \begin{tcolorbox}[enhanced,colback=#1,colframe=popink,boxrule=2pt,arc=9pt,shadow={2.5pt}{-2.5pt}{0pt}{popink},
      left=6pt,right=6pt,top=9pt,bottom=9pt,halign=center,valign=center,height=1.75cm,width=\linewidth,nobeforeafter]
    {\popdisplay\fontsize{12.5}{13}\selectfont\color{white}\MakeUppercase{#2}}\par\vspace{3pt}
    {\centering\popsemi\fontsize{7.8}{9.5}\selectfont\color{white}#3\par}
  \end{tcolorbox}}
\newcommand{\popcontenu}{%
  \par\noindent{\popxboldls\fontsize{7.5}{7.5}\selectfont\color{popmuted}THIS SIGNED COURSE CONTAINS}\par\vspace{7pt}
  \noindent\begin{minipage}[t]{0.235\linewidth}\poptile{popblue}{Course}{complete, illustrated, step by step}\end{minipage}\hfill
  \begin{minipage}[t]{0.235\linewidth}\poptile{poporange}{Methods}{and worked examples}\end{minipage}\hfill
  \begin{minipage}[t]{0.235\linewidth}\poptile{poppink}{Exercises}{exam style + solutions}\end{minipage}\hfill
  \begin{minipage}[t]{0.235\linewidth}\poptile{popgreen}{Summary sheet}{the essentials to remember}\end{minipage}\par}

% Sommaire
\newtcolorbox{sommaire}{enhanced,colback=white,colframe=popink,boxrule=2.2pt,arc=10pt,
  drop shadow=popink,left=18pt,right=18pt,top=13pt,bottom=13pt,before skip=6pt,after skip=20pt,
  before upper={\poppill{Contents}{poppurple}{white}\par\vspace{9pt}\popsemi\fontsize{10.6}{14.5}\selectfont\color{popink}}}

% Signature de fin de cours — poussée tout en bas de la dernière page
\newcommand{\findecours}{%
  \par\vfill\nopagebreak
  \begin{center}\popsquiggle{poporange}\end{center}\vspace{4pt}
  \signatureonbuch
}
\AtBeginDocument{\def\llbracket{\mathopen{[\mkern-3.5mu[}}\def\rrbracket{\mathclose{]\mkern-3.5mu]}}} % intervalles d'entiers (glyphe ⟦ absent de la police)
% Glyphes absents de Fira Math : redéfinis à partir de glyphes disponibles
\AtBeginDocument{%
  \def\setminus{\mathbin{\backslash}}\let\smallsetminus\setminus
  \def\vdots{\mathord{\vbox{\baselineskip3.2pt\lineskiplimit0pt\kern4pt\hbox{.}\hbox{.}\hbox{.}}}}%
  \def\ddots{\mathinner{\mkern1mu\raise6.5pt\hbox{.}\mkern2mu\raise3.6pt\hbox{.}\mkern2mu\raise0.7pt\hbox{.}\mkern1mu}}%
  \def\triangle{\mathord{\scalebox{0.9}{$\symup{\Delta}$}}}%
  \def\nabla{\mathord{\raisebox{\depth}{\scalebox{1}[-1]{$\symup{\Delta}$}}}}%
  \def\checkmark{\popcheck{popdark}}%
  \def\star{\mathbin{\ast}}\def\bigstar{\mathord{\ast}}%
  \def\diamond{\mathbin{\scalebox{0.6}{$\Diamond$}}}%
  \def\bigcup{\mathop{\mathchoice{\raisebox{-0.3em}{\scalebox{1.7}{$\cup$}}}{\scalebox{1.25}{$\cup$}}{\cup}{\cup}}\displaylimits}%
  \def\bigcap{\mathop{\mathchoice{\raisebox{-0.3em}{\scalebox{1.7}{$\cap$}}}{\scalebox{1.25}{$\cap$}}{\cap}{\cap}}\displaylimits}%
  \def\lesseqqgtr{\mathrel{\lessgtr}}\def\bigoplus{\mathop{\oplus}\displaylimits}%
}
\providecommand{\ssi}{if and only if} % abréviation fréquente
\AtBeginDocument{\providecommand{\wideparen}{\overparen}} % arc de cercle

%% FIGFIX-BEGIN v3 — correctifs globaux des figures (docs/onbuch-generation/tools/figfix)
\usepackage{adjustbox}\usepackage{etoolbox}
% toute figure TikZ/pgfplots plus large que la ligne est réduite à la largeur disponible
\BeforeBeginEnvironment{tikzpicture}{\begin{adjustbox}{max width=\linewidth}}
\AfterEndEnvironment{tikzpicture}{\end{adjustbox}}
% légendes pgfplots lisibles par-dessus les courbes
\pgfplotsset{every axis/.append style={legend style={fill=white,fill opacity=0.92,text opacity=1,draw=black!15,inner sep=3pt}}}
% étiquettes d'arêtes (syntaxe edge["..."]) avec fond blanc
\tikzset{every edge quotes/.append style={fill=white,inner sep=1.2pt}}
% bulles de cartes mentales : texte à la ligne dans la bulle, bulle un peu plus grande
\tikzset{every concept/.append style={text width=2.0cm,align=center,minimum size=2.3cm,inner sep=2pt}}
%% FIGFIX-END
\begin{document}

\pagegarde{Support \& locomotion}{Biology}{Upper Sixth}{Upper Sixth — Biology}{5}{4 periods}{This chapter studies the biological basis of support and locomotion in animals: the types of skeletons, the structure and functions of the vertebrate skeleton, the joints and the mechanisms ensuring smooth movement, and the molecular mechanism of muscle contraction. It closes with the concepts of sustainable development, linking biological resource use to the future of Cameroonian society.
\vspace{4pt}
\begin{multicols}{2}\small
\begin{itemize}
\item By the end of this chapter I can define a skeleton and compare hydrostatic, exoskeleton and endoskeleton types with examples.
\item By the end of this chapter I can describe the general plan of the vertebrate skeleton (axial and appendicular) and label its main bones.
\item By the end of this chapter I can state the functions of the skeleton: support, protection, movement, blood cell production and mineral storage.
\item By the end of this chapter I can classify joints (fixed, slightly movable, freely movable/synovial) and describe the structure of a synovial joint.
\item By the end of this chapter I can explain how cartilage, synovial fluid, ligaments and antagonistic muscles ensure smooth movement at joints.
\item By the end of this chapter I can describe the structure of a skeletal muscle from the whole muscle down to the sarcomere.
\end{itemize}\end{multicols}}

\begin{sommaire}
\begin{multicols}{2}
\begin{enumerate}
\item Skeletons and the vertebrate skeleton
\item Joints and smooth movement
\item Muscle contraction
\item Concepts of sustainable development
\item Integration activity
\item Methods and worked examples
\item Exercises
\item Solutions to the exercises
\item Summary sheet
\end{enumerate}
\end{multicols}
\end{sommaire}

\begin{objectifs}
\begin{itemize}
\item By the end of this chapter I can define a skeleton and compare hydrostatic, exoskeleton and endoskeleton types with examples.
\item By the end of this chapter I can describe the general plan of the vertebrate skeleton (axial and appendicular) and label its main bones.
\item By the end of this chapter I can state the functions of the skeleton: support, protection, movement, blood cell production and mineral storage.
\item By the end of this chapter I can classify joints (fixed, slightly movable, freely movable/synovial) and describe the structure of a synovial joint.
\item By the end of this chapter I can explain how cartilage, synovial fluid, ligaments and antagonistic muscles ensure smooth movement at joints.
\item By the end of this chapter I can describe the structure of a skeletal muscle from the whole muscle down to the sarcomere.
\item By the end of this chapter I can explain the sliding filament theory of muscle contraction, including the roles of actin, myosin, ATP, calcium ions and troponin-tropomyosin.
\item By the end of this chapter I can relate muscle contraction to energy supply (aerobic respiration, anaerobic respiration, creatine phosphate) and to fatigue.
\item By the end of this chapter I can define sustainable development and explain its three pillars (environmental, economic, social) with Cameroonian examples.
\item By the end of this chapter I can analyse how biological principles of resource use inform sustainable practices in agriculture, forestry and health.
\end{itemize}
\end{objectifs}

\begin{prerequis}
\begin{itemize}
\item Cell structure and the role of ATP in cellular energy (Lower Sixth)
\item Aerobic and anaerobic respiration and their energy yields (Lower Sixth)
\item Protein structure and the properties of fibrous proteins (Lower Sixth)
\item Basic organisation of the human body: tissues, organs, systems (Lower Sixth)
\item Nervous coordination and the reflex arc (Upper Sixth, coordination chapter)
\end{itemize}
\end{prerequis}

\newpage

\coursec{Skeletons and the vertebrate skeleton}

At the Mokolo market in Yaound\'e, a porter carries a 60 kg bag of cocoa on his head for hours without his spine collapsing, while a footballer at the Ahmadou Ahidjo stadium tears a ligament in a single twisting movement. Why does the same skeleton resist enormous loads in one case and fail in another? How do the muscles of a fisherman paddling on the Wouri river contract thousands of times without rest? And as Cameroon debates the exploitation of its forests and the health of its growing population, how do living systems --- including our own bodies --- teach us about using resources sustainably? This chapter answers these questions by exploring how skeletons support, how joints move, how muscles contract, and how the concept of sustainable development applies to life and society.

In this first section we ask the most basic question: \emph{what is a skeleton, and why do animals need one?}

\courssub{What is a skeleton?}

\textbf{Observation.} Pick up a freshly killed fish and it flops in your hand; yet in water the same fish darts with precision. Press the body of an earthworm and it is soft, yet it can push through compacted soil. Hold a crab shell and it is hard as a nut, yet the animal inside moved nimbly. In every case, something inside or around the body resists deformation and gives muscles something to pull against. That ``something'' is the \cle{skeleton}.

\begin{definition}[Skeleton]
A \cle{skeleton} is a rigid or semi-rigid framework that \textbf{supports} the body and maintains its shape, \textbf{protects} delicate internal organs, and provides \textbf{attachment surfaces (leverage)} for muscles so that movement is possible.
\end{definition}

The three roles are worth memorising as a triad, because examiners love to ask for them:

\begin{enumerate}
\item \textbf{Support:} the skeleton bears the weight of the body against gravity (the porter's vertebral column under the cocoa bag) and maintains body shape.
\item \textbf{Protection:} hard parts shield soft organs --- the skull protects the brain, the rib cage protects the heart and lungs.
\item \textbf{Leverage for movement:} muscles can only \emph{pull}, never push. They pull on skeletal elements which act as levers, producing movement at joints.
\end{enumerate}

Three great engineering solutions to the support problem have evolved in animals: the \cle{hydrostatic skeleton}, the \cle{exoskeleton} and the \cle{endoskeleton}.

\courssub{Hydrostatic skeletons}

\textbf{Observation.} An earthworm has no bone and no shell, yet it is firm to the touch and can elongate, shorten and burrow. Squeeze a sealed plastic bottle full of water: it resists, because water is \emph{incompressible}. The earthworm uses exactly this principle.

\begin{definition}[Hydrostatic skeleton]
A \cle{hydrostatic skeleton} consists of a fluid-filled body cavity surrounded by a muscular wall. Because the fluid is incompressible, contraction of the muscles changes the \emph{shape} of the animal but not its \emph{volume}; the fluid transmits the pressure and stiffens the body, providing a firm structure against which muscles act.
\end{definition}

In the earthworm, the body cavity (coelom) is divided into segments, each a separate fluid compartment. Circular muscles make a segment long and thin; longitudinal muscles make it short and fat. Alternating these contractions along the body, with bristles (chaetae) gripping the soil, produces burrowing locomotion. The sea anemone uses the same trick: it fills its gastrovascular cavity with water and closes its mouth, turning its soft body into a stiff column.

\begin{attention}[Limits of hydrostatic skeletons]
A hydrostatic skeleton offers almost no \textbf{protection} against predators or blows, and it works poorly for large animals on land, where gravity demands a rigid framework. This is why all large terrestrial animals have hard skeletons.
\end{attention}

\courssub{Exoskeletons}

\textbf{Observation.} Crack a crab claw at the seaside in Kribi: the animal's ``skeleton'' is on the \emph{outside}. Arthropods (crabs, insects, spiders, millipedes) secrete a tough external \cle{cuticle} containing \cle{chitin}, a strong, light polysaccharide, often reinforced with calcium salts (in crustaceans) or hardened proteins (in insects).

\begin{definition}[Exoskeleton]
An \cle{exoskeleton} is a hard, non-living external covering secreted by the epidermis. It supports and protects the body and provides attachment points for muscles on its inner surface.
\end{definition}

\textbf{Advantages:}
\begin{itemize}
\item Excellent \textbf{protection} against predators, desiccation and mechanical damage --- a major reason insects conquered land.
\item Muscles attach directly to the inside of rigid plates, giving efficient leverage.
\item Joints between plates allow precise, varied movement.
\end{itemize}

\textbf{Limits:}
\begin{itemize}
\item The cuticle cannot grow. The animal must periodically shed it --- \cle{moulting} or \cle{ecdysis} --- and secrete a new, larger one. During moulting the animal is soft and vulnerable.
\item An external skeleton heavy enough to support a very large body would be impossibly thick; this sets an upper size limit on arthropods.
\end{itemize}

\courssub{Endoskeletons}

Vertebrates --- fish, amphibians, reptiles, birds and mammals --- took the opposite route: an \emph{internal} framework of \cle{cartilage} and \cle{bone}.

\begin{definition}[Endoskeleton]
An \cle{endoskeleton} is an internal living framework of cartilage and/or bone. Because it is living tissue, it \textbf{grows with the body}, so no moulting is needed, and it can repair itself after damage.
\end{definition}

\begin{propriete}[Bone is a living tissue]
Bone is not dead scaffolding. It is penetrated by \textbf{blood vessels} and \textbf{nerves}, contains living cells (osteoblasts which build bone, osteoclasts which resorb it), \textbf{grows} during development and can \textbf{repair itself} after a fracture. (Proof not examinable; the property itself is.)
\end{propriete}

\textbf{Bone versus cartilage.} Both are connective tissues built by cells embedded in a matrix they secrete, but they differ decisively:

\begin{center}
\begin{tabularx}{\linewidth}{|l|X|X|}
\hline
\rowcolor{popblueL} \textbf{Feature} & \textbf{Cartilage} & \textbf{Bone} \\
\hline
Matrix & Firm gel of collagen fibres in a rubbery ground substance & Collagen fibres hardened by \textbf{calcium phosphate} (hydroxyapatite) crystals \\
\hline
Blood vessels & \textbf{None} --- cells nourished by diffusion & Richly supplied with blood and nerves \\
\hline
Rigidity & Flexible, semi-rigid & Hard and rigid \\
\hline
Growth/repair & Slow, poor repair & Grows and repairs well \\
\hline
Where found & Nose, ear, trachea rings, joint surfaces, embryonic skeleton & Adult skeleton of most vertebrates \\
\hline
\end{tabularx}
\end{center}

\begin{attention}[Classic confusion]
Do not say ``cartilage is soft bone.'' Cartilage contains \textbf{no calcium salts} and \textbf{no blood vessels}; bone contains both. In the embryo the skeleton is first laid down in cartilage, which is then progressively replaced by bone --- a process called \cle{ossification}. In outline: cartilage cells die, blood vessels invade, osteoblasts deposit calcium phosphate on the collagen framework, and hard bone results. Ossification continues into the late teens, which is why children's bones are more flexible.
\end{attention}

\begin{popfigure}
\begin{center}
\begin{tikzpicture}[scale=0.9, font=\small]
% Earthworm - hydrostatic
\draw[fill=popgreenL, draw=popgreen, thick] (0,0) ellipse (2.2 and 0.7);
\draw[popgreen, thick] (0,0) ellipse (1.6 and 0.4);
\fill[popblue!40] (0,0) ellipse (1.55 and 0.35);
\node[align=center] at (0,-1.3) {\textbf{Earthworm}\\ hydrostatic};
\node[align=center, text=popink, text width=2.5cm] at (3.4,0.5) {fluid-filled cavity};
\draw[->, popink] (2.6,0.35) -- (1.4,0.15);
\node[align=center, text=popink, text width=2.5cm] at (3.4,-0.6) {muscular wall};
\draw[->, popink] (2.7,-0.45) -- (1.9,-0.25);
% Crab - exoskeleton
\draw[fill=poporange!35, draw=poporange, very thick] (8,0.4) ellipse (1.5 and 0.8);
\draw[poporange, very thick] (6.9,1.0) -- (6.3,1.7);
\draw[poporange, very thick] (9.1,1.0) -- (9.7,1.7);
\draw[poporange, very thick] (6.6,0.6) -- (5.9,1.0);
\draw[poporange, very thick] (9.4,0.6) -- (10.1,1.0);
\fill[white] (7.6,0.6) circle (0.08); \fill[white] (8.4,0.6) circle (0.08);
\node[align=center] at (8,-1.3) {\textbf{Crab}\\ exoskeleton (chitin)};
\node[align=center, text=popink, text width=2.5cm] at (11.2,0.9) {hard cuticle outside};
\draw[->, popink] (10.3,0.75) -- (9.2,0.7);
% Fish - endoskeleton
\draw[fill=popblueL, draw=popblue, thick] (14.2,0.2) .. controls (16.2,1.0) and (17.8,0.9) .. (18.6,0.2) .. controls (17.8,-0.5) and (16.2,-0.6) .. (14.2,0.2);
\draw[fill=popblue, draw=popblue] (18.6,0.2) -- (19.5,0.8) -- (19.5,-0.4) -- cycle;
\draw[popdark, thick] (14.5,0.2) -- (18.4,0.2);
\foreach \x in {14.9,15.4,15.9,16.4,16.9,17.4,17.9} {
  \draw[popdark] (\x,0.2) -- (\x,0.55);
  \draw[popdark] (\x,0.2) -- (\x,-0.15);
}
\fill[black] (14.7,0.4) circle (0.05);
\node[align=center] at (16.5,-1.3) {\textbf{Fish}\\ endoskeleton};
\node[align=center, text=popink, text width=2.5cm] at (16.5,1.2) {internal backbone};
\draw[->, popink] (16.5,0.95) -- (16.5,0.35);
\end{tikzpicture}
\end{center}
\legende{The three types of skeleton: hydrostatic (earthworm), exoskeleton (crab) and endoskeleton (fish).}
\end{popfigure}

\courssub{The general plan of the vertebrate skeleton}

The mammalian skeleton is organised in two great divisions:

\begin{definition}[Axial and appendicular skeleton]
The \cle{axial skeleton} forms the central axis of the body: the \textbf{skull}, the \textbf{vertebral column}, the \textbf{ribs} and the \textbf{sternum}. The \cle{appendicular skeleton} comprises the \textbf{limbs} and the \textbf{girdles} (pectoral and pelvic) that attach them to the axis.
\end{definition}

\begin{methode}[Distinguishing axial from appendicular on any diagram]
\begin{enumerate}
\item Find the \textbf{vertebral column} --- the chain of vertebrae down the midline. Everything continuous with it (skull above, ribs and sternum branching off) is \textbf{axial}.
\item Find the \textbf{girdles}: the shoulder (pectoral) girdle --- scapula and clavicle --- and the hip (pelvic) girdle. These, plus every bone of the limbs attached to them, are \textbf{appendicular}.
\item Rule of thumb: \emph{axis = midline and rib cage; appendicular = girdles + limbs.}
\end{enumerate}
\end{methode}

\begin{popfigure}
\begin{center}
\begin{tikzpicture}[scale=0.62, font=\scriptsize]
% ---- AXIAL (popblue) ----
% skull
\draw[fill=popblueL, draw=popblue, thick] (0,8.6) circle (0.75);
\draw[fill=popblueL, draw=popblue, thick] (0,7.85) ellipse (0.4 and 0.28);
% vertebral column
\foreach \y in {7.4,7.0,6.6,6.2,5.8,5.4,5.0,4.6,4.2,3.8} {
  \draw[fill=popblueL, draw=popblue] (-0.22,\y) rectangle (0.22,\y+0.28);
}
% ribs
\foreach \y/\w in {6.9/1.15, 6.5/1.35, 6.1/1.45, 5.7/1.35, 5.3/1.1} {
  \draw[popblue, thick] (0,\y) .. controls (\w,\y+0.15) and (\w+0.25,\y-0.25) .. (\w-0.15,\y-0.55);
  \draw[popblue, thick] (0,\y) .. controls (-\w,\y+0.15) and (-\w-0.25,\y-0.25) .. (-\w+0.15,\y-0.55);
}
% sternum
\draw[popblue, very thick] (0,6.9) -- (0,5.0);
% ---- APPENDICULAR (poporange) ----
% pectoral girdle
\draw[poporange, very thick] (-0.2,7.3) -- (-1.5,7.1);
\draw[poporange, very thick] (0.2,7.3) -- (1.5,7.1);
\draw[fill=poporange!30, draw=poporange] (-1.7,6.9) ellipse (0.3 and 0.45);
\draw[fill=poporange!30, draw=poporange] (1.7,6.9) ellipse (0.3 and 0.45);
% arms
\draw[poporange, very thick] (-1.7,6.5) -- (-2.1,5.2);
\draw[poporange, very thick] (1.7,6.5) -- (2.1,5.2);
\draw[poporange, thick] (-2.1,5.2) -- (-2.35,3.9);
\draw[poporange, thick] (-2.05,5.15) -- (-2.15,3.9);
\draw[poporange, thick] (2.1,5.2) -- (2.35,3.9);
\draw[poporange, thick] (2.05,5.15) -- (2.15,3.9);
\draw[poporange, thick] (-2.25,3.9) circle (0.16);
\draw[poporange, thick] (2.25,3.9) circle (0.16);
% pelvic girdle
\draw[fill=poporange!30, draw=poporange, thick] (-0.9,3.7) ellipse (0.75 and 0.55);
\draw[fill=poporange!30, draw=poporange, thick] (0.9,3.7) ellipse (0.75 and 0.55);
% legs
\draw[poporange, very thick] (-0.9,3.2) -- (-1.0,1.9);
\draw[poporange, very thick] (0.9,3.2) -- (1.0,1.9);
\draw[poporange, thick] (-1.0,1.9) -- (-1.0,0.6);
\draw[poporange, thick] (-0.92,1.9) -- (-0.85,0.6);
\draw[poporange, thick] (1.0,1.9) -- (1.0,0.6);
\draw[poporange, thick] (0.92,1.9) -- (0.85,0.6);
\draw[poporange, thick] (-1.0,0.6) -- (-0.6,0.45);
\draw[poporange, thick] (1.0,0.6) -- (1.35,0.45);
% labels
\node[text=popblue, align=center] at (-3.6,8.6) {skull\\ (axial)};
\draw[->, popblue] (-2.6,8.6) -- (-0.85,8.6);
\node[text=popblue, align=center] at (3.6,6.4) {ribs +\\ sternum};
\draw[->, popblue] (2.7,6.3) -- (1.4,6.2);
\node[text=popblue, align=center] at (-3.4,5.6) {vertebral\\ column};
\draw[->, popblue] (-2.4,5.5) -- (-0.3,5.5);
\node[text=poporange, align=center] at (4.0,7.3) {pectoral\\ girdle};
\draw[->, poporange] (3.1,7.2) -- (1.9,7.0);
\node[text=poporange, align=center] at (-4.0,3.5) {pelvic\\ girdle};
\draw[->, poporange] (-3.0,3.5) -- (-1.7,3.6);
\node[text=poporange, align=center] at (3.9,1.4) {limbs\\ (appendicular)};
\draw[->, poporange] (3.0,1.3) -- (1.15,1.3);
\end{tikzpicture}
\end{center}
\legende{The human skeleton: axial skeleton in blue (skull, vertebral column, ribs, sternum), appendicular skeleton in orange (girdles and limbs).}
\end{popfigure}

\courssub{Regional specialisation of the vertebrae}

The vertebral column is not a uniform rod: its vertebrae are specialised region by region, each form matching its function --- a favourite GCE theme (``structure related to function'').

\begin{center}
\begin{tabularx}{\linewidth}{|l|l|X|}
\hline
\rowcolor{popblueL} \textbf{Region} & \textbf{Number (man)} & \textbf{Specialisation and function} \\
\hline
Cervical (neck) & 7 & Small, light; support the head and allow it to turn. \textbf{Atlas} (C1) is a ring that holds the skull and allows nodding; \textbf{axis} (C2) bears a peg (odontoid process) around which the atlas rotates, allowing the head to shake ``no''. \\
\hline
Thoracic (chest) & 12 & Bear facets for articulation with the \textbf{ribs}; long downward neural spines; protect heart and lungs with the rib cage. \\
\hline
Lumbar (lower back) & 5 & Largest, most massive vertebrae; bear the \textbf{weight of the trunk} --- the porter's load passes through them. \\
\hline
Sacral & 5 (fused) & Fused into the \textbf{sacrum}, firmly joined to the pelvic girdle; transmits weight to the legs. \\
\hline
Caudal & 4 (fused) & Fused into the \textbf{coccyx}; vestigial tail. \\
\hline
\end{tabularx}
\end{center}

\courssub{The pentadactyl limb: one plan, many functions}

The limbs of all tetrapod vertebrates are built on the same \cle{pentadactyl} (five-fingered) plan: one proximal bone (humerus/femur), two middle bones (radius and ulna / tibia and fibula), a group of small bones (carpals/tarsals), then five digit rays. Yet look at the outcomes:

\begin{itemize}
\item \textbf{Arm of man:} five free digits for grasping and manipulation.
\item \textbf{Wing of bird:} digits reduced and fused, the limb elongated to support feathers for flight.
\item \textbf{Flipper of whale:} bones shortened and flattened into a paddle for swimming.
\end{itemize}

Same bones, same relative positions, different shapes and sizes. Such \cle{homologous structures} are powerful \textbf{evidence of common ancestry}: the most economical explanation is that all these vertebrates inherited the same limb plan from a shared ancestor, then modified it by evolution.

\begin{savaistu}[A favourite GCE question]
``The pentadactyl limb provides evidence for evolution. Discuss.'' This is one of the most frequently set structured questions on evidence for evolution. Be ready to name the bones (humerus, radius, ulna, carpals, metacarpals, phalanges), to give three modified examples (man, bird, whale --- or bat, horse), and to state the conclusion: homology implies descent with modification from a common ancestor.
\end{savaistu}

\courssub{Structure of a long bone}

A long bone such as the femur or humerus shows the economy of skeletal design:

\begin{popfigure}
\begin{center}
\begin{tikzpicture}[scale=0.95, font=\scriptsize]
% long bone outline
\draw[thick, popdark, fill=popgold!15]
  (-0.9,3.4) .. controls (-1.5,3.9) and (-1.5,4.6) .. (-0.9,5.0)
  .. controls (-0.4,5.35) and (0.4,5.35) .. (0.9,5.0)
  .. controls (1.5,4.6) and (1.5,3.9) .. (0.9,3.4)
  -- (0.55,3.2) -- (0.55,0.2) -- (0.9,0.0)
  .. controls (1.5,-0.5) and (1.5,-1.2) .. (0.9,-1.6)
  .. controls (0.4,-1.95) and (-0.4,-1.95) .. (-0.9,-1.6)
  .. controls (-1.5,-1.2) and (-1.5,-0.5) .. (-0.9,0.0)
  -- (-0.55,0.2) -- (-0.55,3.2) -- cycle;
% marrow cavity
\draw[fill=white, draw=popdark] (-0.3,0.4) rectangle (0.3,3.0);
% spongy bone at ends
\foreach \p in {(-0.55,4.2),(-0.2,4.5),(0.3,4.3),(0.0,4.9),(-0.5,4.7),(0.55,4.75),(-0.55,-0.9),(-0.15,-1.1),(0.3,-0.85),(0.0,-1.5),(-0.45,-1.4),(0.5,-1.35)} {
  \fill[poporange!50] \p circle (0.12);
}
% labels
\node[align=center, text=popink, text width=3cm] at (3.4,4.6) {spongy bone (in epiphysis)};
\draw[->, popink] (2.5,4.5) -- (0.7,4.5);
\node[align=center, text=popink, text width=3cm] at (3.4,3.3) {compact bone (dense shaft wall)};
\draw[->, popink] (2.5,3.2) -- (0.58,2.6);
\node[align=center, text=popink, text width=3cm] at (3.4,1.6) {marrow cavity (yellow marrow)};
\draw[->, popink] (2.5,1.6) -- (0.32,1.7);
\node[align=center, text=popink, text width=3cm] at (-3.4,4.2) {epiphysis (head)};
\draw[->, popink] (-2.5,4.2) -- (-1.1,4.3);
\node[align=center, text=popink, text width=3cm] at (-3.4,1.7) {diaphysis (shaft)};
\draw[->, popink] (-2.5,1.7) -- (-0.6,1.7);
\node[align=center, text=popink, text width=3cm] at (-3.4,0.2) {periosteum (outer membrane)};
\draw[->, popink] (-2.5,0.25) -- (-0.6,0.9);
\node[align=center, text=popink, text width=3cm] at (3.4,-0.9) {cartilage cap (joint surface)};
\draw[->, popink] (2.5,-0.9) -- (0.8,-1.35);
\end{tikzpicture}
\end{center}
\legende{Structure of a long bone: epiphyses with spongy bone, shaft (diaphysis) of compact bone around the marrow cavity, wrapped in periosteum.}
\end{popfigure}

\begin{itemize}
\item \textbf{Compact bone} forms the dense, hard wall of the shaft --- maximum strength where bending forces are greatest.
\item \textbf{Spongy bone}, a lattice of bony struts in the heads, is light yet resists compression; its spaces contain \textbf{red marrow}, where blood cells are manufactured.
\item The \textbf{marrow cavity} of the shaft contains yellow marrow (fat store).
\item The \textbf{periosteum}, a tough membrane covering the bone, carries blood vessels and nerves and is essential for growth in thickness and for fracture repair.
\end{itemize}

\courssub{Functions of the mammalian skeleton --- summary}

\begin{aretenir}[Five functions to know by heart]
\begin{enumerate}
\item \textbf{Support:} maintains body shape and bears weight against gravity.
\item \textbf{Protection:} skull protects the brain, vertebrae protect the spinal cord, ribs and sternum protect the heart and lungs.
\item \textbf{Movement:} bones act as levers pulled by muscles across joints.
\item \textbf{Blood cell manufacture:} red bone marrow produces red blood cells, white blood cells and platelets.
\item \textbf{Mineral reservoir:} bone stores \textbf{calcium} and \textbf{phosphorus}, releasing them into the blood when needed.
\end{enumerate}
\end{aretenir}

\begin{exoresolu}[Worked example --- skeleton types and functions]
\textbf{Statement.} (a) Name the type of skeleton found in (i) an earthworm, (ii) a grasshopper, (iii) a goat. (b) Explain why a grasshopper must moult whereas a goat does not. (c) State two functions of the goat's skeleton other than support.
\tcblower
\textbf{Solution.}
\begin{enumerate}
\item (a) (i) Earthworm: \textbf{hydrostatic skeleton} (fluid-filled coelom). (ii) Grasshopper: \textbf{exoskeleton} of chitin. (iii) Goat: \textbf{endoskeleton} of bone and cartilage.
\item (b) The grasshopper's exoskeleton is a non-living cuticle which cannot grow, so it must be shed (ecdysis) and replaced by a larger one. The goat's endoskeleton is living tissue which grows with the body, so no moulting is necessary.
\item (c) Any two of: protection of delicate organs (brain, spinal cord, heart, lungs); leverage/attachment for muscles enabling movement; manufacture of blood cells in red bone marrow; storage/reservoir of calcium and phosphorus.
\end{enumerate}
\end{exoresolu}

\begin{exercice}[Check your understanding]
\begin{enumerate}
\item Define the term \emph{skeleton} and state its three general roles.
\item Explain the principle of the hydrostatic skeleton and give one named example.
\item Give two advantages and two limitations of the arthropod exoskeleton.
\item Construct a table of three differences between bone and cartilage.
\item On a sketch of any mammal, shade the axial skeleton in one colour and the appendicular skeleton in another, and label the pectoral and pelvic girdles.
\item Explain how the structure of (i) a lumbar vertebra and (ii) the atlas vertebra is related to its function.
\item The wing of a bird and the arm of a man contain the same basic set of bones. What is this evidence of, and what name is given to such structures?
\end{enumerate}
\end{exercice}

\begin{corrige}[Exercise 1]
\begin{enumerate}
\item A skeleton is a rigid or semi-rigid framework that supports the body, protects internal organs and provides attachment/leverage for muscles. Roles: support, protection, movement (leverage).
\item A fluid-filled cavity surrounded by muscle; the fluid is incompressible, so muscle contraction changes shape but not volume, stiffening the body as a firm base for movement. Example: earthworm (or sea anemone).
\item Advantages: protection (against predators/desiccation); firm muscle attachment. Limitations: must be moulted for growth (animal vulnerable during ecdysis); limits maximum body size.
\item Bone: matrix hardened with calcium phosphate; has blood vessels and nerves; repairs well. Cartilage: no calcium salts; no blood vessels; flexible, poor repair.
\item Axial: skull, vertebral column, ribs, sternum. Appendicular: scapula + clavicle (pectoral girdle), pelvic girdle, all limb bones.
\item (i) Lumbar vertebrae are large and massive to bear the weight of the trunk. (ii) The atlas is a ring supporting the skull and pivoting on the axis to allow head rotation.
\item Evidence of common ancestry (evolution); the structures are called homologous structures (the pentadactyl limb plan).
\end{enumerate}
\end{corrige}

\coursec{Joints and smooth movement}

In the previous section we studied the skeleton as a framework of rigid bones. But a pile of rigid bones cannot move: a skeleton made of a single solid piece would be as useless for locomotion as a statue. Movement is only possible because bones meet one another at \cle{joints}, and because muscles pull across those joints. Watch a student at Lyc\'ee de Nkolbisson writing, a footballer in Douala kicking a ball, or a farmer in the North West bending to plant maize: every one of these actions depends on joints working smoothly. In this section we define joints, classify them, study the structure of synovial joints in detail, and explain how antagonistic muscles and bony levers produce controlled movement.

\courssub{Definition and classification of joints}

\begin{definition}[Joint (articulation)]
A \cle{joint} or \cle{articulation} is a point where two or more bones meet. Joints hold the bones of the skeleton together and, in most cases, allow movement between them, so that the skeleton can act as a system of levers for locomotion.
\end{definition}

Joints are classified according to the \textbf{degree of movement} they permit.

\begin{center}
\begin{tabularx}{\linewidth}{|l|X|X|}
\hline
\rowcolor{popblueL} \textbf{Class} & \textbf{Description} & \textbf{Examples} \\
\hline
Fixed (immovable) & Bones tightly locked together; no movement possible & Sutures between the bones of the cranium (skull) \\
\hline
Slightly movable & Bones joined by cartilage or tough fibrous tissue; limited movement & Between adjacent vertebrae; pubic symphysis (between the two pelvic bones) \\
\hline
Freely movable (synovial) & Bone ends separated by a fluid-filled cavity; wide range of movement & Shoulder, hip, elbow, wrist \\
\hline
\end{tabularx}
\end{center}

\textbf{Fixed joints} protect delicate organs: the sutures of the skull form a rigid case around the brain. \textbf{Slightly movable joints} combine strength with a little flexibility: the pads of cartilage (intervertebral discs) between vertebrae let the backbone bend and also act as shock absorbers when we walk or jump. The pubic symphysis softens slightly during pregnancy, helping childbirth. \textbf{Freely movable joints} are the joints of locomotion proper, and we now study them in detail.

\courssub{Types of synovial joints}

\begin{definition}[Synovial joint]
A \cle{synovial joint} is a freely movable joint in which the articulating ends of the bones are separated by a fluid-filled cavity, the \cle{synovial cavity}.
\end{definition}

Synovial joints are subdivided according to the \textbf{shape of the articulating surfaces}, which determines the movements allowed.

\begin{center}
\begin{tabularx}{\linewidth}{|l|X|X|}
\hline
\rowcolor{popblueL} \textbf{Type} & \textbf{Movement allowed} & \textbf{Example} \\
\hline
Hinge & Movement in one plane only (flexion and extension), like a door hinge & Elbow, knee \\
\hline
Ball-and-socket & Movement in all planes, including rotation; the widest range & Hip, shoulder \\
\hline
Pivot & Rotation of one bone around another & Atlas on axis (neck), allowing the head to turn \\
\hline
Gliding & Flat surfaces slide over each other; small movements & Between the small bones of the wrist and ankle \\
\hline
Saddle & Movement in two planes; each surface is concave in one direction and convex in the other & Base of the thumb (allows the thumb to touch each finger) \\
\hline
\end{tabularx}
\end{center}

\begin{methode}[Identifying the type of synovial joint]
To identify a synovial joint in an examination question, proceed in order:
\begin{enumerate}
\item Look at the \textbf{shape of the articulating surfaces}: a rounded head in a cup indicates ball-and-socket; a cylinder in a ring indicates pivot; flat surfaces indicate gliding.
\item Ask \textbf{how many planes of movement} are possible: one plane suggests hinge; two planes suggest saddle; all planes suggest ball-and-socket.
\item Use the \textbf{position in the body} to confirm: joints needing great mobility (shoulder, hip) are ball-and-socket; joints needing strength in one direction (elbow, knee) are hinge joints.
\end{enumerate}
\end{methode}

\courssub{Structure of a synovial joint}

\begin{popfigure}
\begin{tikzpicture}[scale=1.0]
% upper bone
\fill[popblueL] (-2.2,1.2) rectangle (2.2,3.2);
\draw[thick,popdark] (-2.2,1.2) -- (-2.2,3.2) -- (2.2,3.2) -- (2.2,1.2);
\draw[thick,popdark] (-2.2,1.2) .. controls (-1.2,0.7) and (1.2,0.7) .. (2.2,1.2);
% lower bone
\fill[popblueL] (-2.2,-3.2) rectangle (2.2,-1.2);
\draw[thick,popdark] (-2.2,-1.2) -- (-2.2,-3.2) -- (2.2,-3.2) -- (2.2,-1.2);
\draw[thick,popdark] (-2.2,-1.2) .. controls (-1.2,-0.7) and (1.2,-0.7) .. (2.2,-1.2);
% cartilage
\draw[very thick,popgreen] (-2.2,1.2) .. controls (-1.2,0.7) and (1.2,0.7) .. (2.2,1.2);
\draw[very thick,popgreen] (-2.2,-1.2) .. controls (-1.2,-0.7) and (1.2,-0.7) .. (2.2,-1.2);
% synovial cavity
\fill[popgold!40] (-1.9,0.55) .. controls (-1.0,0.1) and (1.0,0.1) .. (1.9,0.55)
 -- (1.9,-0.55) .. controls (1.0,-0.1) and (-1.0,-0.1) .. (-1.9,-0.55) -- cycle;
% capsule
\draw[very thick,poporange] (-2.2,1.6) -- (-2.9,1.6) -- (-2.9,-1.6) -- (-2.2,-1.6);
\draw[very thick,poporange] (2.2,1.6) -- (2.9,1.6) -- (2.9,-1.6) -- (2.2,-1.6);
% synovial membrane (inner line of capsule)
\draw[thick,poppink,dashed] (-2.75,1.4) -- (-2.75,-1.4);
\draw[thick,poppink,dashed] (2.75,1.4) -- (2.75,-1.4);
% ligaments
\draw[line width=2.5pt,popink] (-3.4,1.9) -- (-3.4,-1.9);
\draw[line width=2.5pt,popink] (3.4,1.9) -- (3.4,-1.9);
% labels
\node[popdark] at (0,2.3) {\small Bone};
\node[popdark] at (0,-2.3) {\small Bone};
\draw[->,popgreen] (4.6,1.1) -- (1.4,0.95);
\node[popgreen,anchor=west] at (4.6,1.1) {\small Articular cartilage};
\draw[->,popdark] (4.6,0.1) -- (1.0,0.05);
\node[popdark,anchor=west] at (4.6,0.1) {\small Synovial fluid in cavity};
\draw[->,poppink] (-4.8,0.9) -- (-2.78,0.6);
\node[poppink,anchor=east,align=right] at (-4.8,0.9) {\small Synovial\\[-2pt]\small membrane};
\draw[->,poporange] (-4.8,-0.3) -- (-2.9,-0.3);
\node[poporange,anchor=east] at (-4.8,-0.3) {\small Joint capsule};
\draw[->,popink] (-4.8,-1.4) -- (-3.4,-1.2);
\node[popink,anchor=east] at (-4.8,-1.4) {\small Ligament};
\end{tikzpicture}
\legende{Longitudinal section through a synovial joint (e.g. the knee), showing the main structures.}
\end{popfigure}

The structures of a synovial joint, and their functions, are:

\begin{itemize}
\item \textbf{Articular cartilage}: a smooth, slippery layer of cartilage covering the ends of the bones. It \cle{reduces friction} between the bones and \cle{absorbs shock} when the joint is jarred (for example when landing from a jump).
\item \textbf{Synovial membrane}: the inner lining of the joint capsule; it \cle{secretes synovial fluid}.
\item \textbf{Synovial fluid}: an oily fluid filling the synovial cavity.
\item \textbf{Joint capsule}: a tough fibrous sleeve enclosing the joint and holding the synovial fluid in place.
\item \textbf{Ligaments}: strong, slightly elastic bands of fibrous tissue that \cle{join bone to bone} across the joint, holding the bones together while still allowing movement.
\item \textbf{Tendons}: tough, inelastic cords that \cle{join muscle to bone}, transmitting the pull of muscles across the joint.
\end{itemize}

\begin{propriete}[Double role of synovial fluid]
Synovial fluid has two essential functions: it \cle{lubricates} the joint so that the cartilage surfaces slide over each other with very little friction, and it \cle{supplies nutrients} (and removes wastes) for the articular cartilage, which is \emph{avascular} --- it contains no blood vessels and must be fed by diffusion from the fluid.
\end{propriete}

\begin{attention}[Ligaments versus tendons]
A very common examination mistake is to confuse these two structures. Remember: \textbf{ligaments join bone to bone} and hold the joint together; \textbf{tendons join muscle to bone} and transmit the force of contraction. Both are made of fibrous connective tissue, but their attachments and roles are different.
\end{attention}

\courssub{Antagonistic muscles and movement at joints}

Muscles can only \textbf{pull}; they cannot push. A single muscle can therefore move a bone in one direction only. To move the bone back, a second muscle pulling the other way is needed.

\begin{definition}[Antagonistic muscles]
\cle{Antagonistic muscles} are pairs of muscles whose contractions produce \emph{opposite} movements at a joint: when one muscle (the \cle{agonist}) contracts, its partner (the \cle{antagonist}) relaxes, and vice versa.
\end{definition}

At the elbow, the \textbf{biceps} (in front of the upper arm) is a \cle{flexor}: its contraction \emph{flexes} (bends) the arm. The \textbf{triceps} (behind the upper arm) is an \cle{extensor}: its contraction \emph{extends} (straightens) the arm. Each muscle is attached to bone at two points: the \cle{origin} is the attachment on the bone that stays still, and the \cle{insertion} is the attachment on the bone that moves. The biceps originates on the scapula and inserts on the radius; when it contracts, the radius (and the forearm) is pulled upwards.

\begin{popfigure}
\begin{tikzpicture}[scale=0.95]
% ===== FLEXION (left) =====
% upper arm bone
\draw[thick,popdark,fill=popblueL] (-6.2,0.4) rectangle (-5.4,3.0);
% forearm bone raised
\draw[thick,popdark,fill=popblueL,rotate around={55:(-5.8,0.4)}] (-6.2,0.4) rectangle (-5.4,-2.4);
% biceps contracted (thick)
\draw[line width=5pt,poporange] (-6.55,2.9) .. controls (-7.3,1.6) .. (-6.9,0.9);
% triceps relaxed (thin)
\draw[line width=1.5pt,popgreen] (-5.05,2.9) -- (-5.05,0.1);
\node[poporange,align=center] at (-8.3,1.9) {\small Biceps\\[-2pt]\small contracted};
\draw[->,poporange] (-7.6,1.7) -- (-7.15,1.6);
\node[popgreen,align=center] at (-3.9,1.9) {\small Triceps\\[-2pt]\small relaxed};
\draw[->,popgreen] (-4.6,1.6) -- (-5.0,1.5);
\node[popdark] at (-5.8,-2.6) {\small \textbf{Flexion}};
% ===== EXTENSION (right) =====
\draw[thick,popdark,fill=popblueL] (1.6,0.4) rectangle (2.4,3.0);
\draw[thick,popdark,fill=popblueL] (1.6,0.4) rectangle (2.4,-2.4);
% biceps relaxed
\draw[line width=1.5pt,poporange] (1.25,2.9) -- (1.25,-1.9);
% triceps contracted
\draw[line width=5pt,popgreen] (2.75,2.9) .. controls (3.5,1.4) .. (2.75,-1.6);
\node[poporange,align=center] at (-0.2,1.9) {\small Biceps\\[-2pt]\small relaxed};
\draw[->,poporange] (0.5,1.6) -- (1.2,1.5);
\node[popgreen,align=center] at (4.6,1.9) {\small Triceps\\[-2pt]\small contracted};
\draw[->,popgreen] (3.9,1.6) -- (3.35,1.5);
\node[popdark] at (2.0,-2.9) {\small \textbf{Extension}};
\end{tikzpicture}
\legende{Antagonistic action of biceps and triceps at the elbow: flexion (left) and extension (right).}
\end{popfigure}

\begin{exemplebox}[Raising and lowering a load]
When a porter at a market in Bamenda lifts a basket onto his head, the biceps contracts (agonist) and the triceps relaxes (antagonist): the elbow flexes. To lower the basket in a controlled way, the triceps contracts while the biceps relaxes. The two muscles never contract hard at the same time in normal movement; their coordination is controlled by the nervous system.
\end{exemplebox}

\courssub{Bones as levers}

When a muscle pulls on a bone across a joint, the bone behaves as a \cle{lever}. Every lever has three elements: the \cle{fulcrum} (the pivot, here the joint), the \cle{effort} (the pull of the muscle) and the \cle{load} (the resistance moved). According to the relative positions of these three elements, three orders (classes) of levers exist.

\begin{popfigure}
\begin{tikzpicture}[scale=0.9]
% Class 1
\draw[thick,popdark] (-6.5,0) -- (-2.5,0);
\fill[poporange] (-4.5,-0.7) -- (-4.9,0) -- (-4.1,0) -- cycle;
\node[poporange] at (-4.5,-1.1) {\small F};
\draw[->,popink,thick] (-6.3,0.9) -- (-6.3,0.1);
\node[popink] at (-6.3,1.2) {\small E};
\draw[->,popgreen,thick] (-2.7,0.9) -- (-2.7,0.1);
\node[popgreen] at (-2.7,1.2) {\small L};
\node[popdark,align=center] at (-4.5,-1.9) {\small Class 1: E--F--L\\[-2pt]\small (nodding the head)};
% Class 2
\draw[thick,popdark] (-0.5,0) -- (3.5,0);
\fill[poporange] (-0.3,-0.7) -- (-0.7,0) -- (0.1,0) -- cycle;
\node[poporange] at (-0.3,-1.1) {\small F};
\draw[->,popgreen,thick] (1.5,0.9) -- (1.5,0.1);
\node[popgreen] at (1.5,1.2) {\small L};
\draw[->,popink,thick] (3.3,0.9) -- (3.3,0.1);
\node[popink] at (3.3,1.2) {\small E};
\node[popdark,align=center] at (1.5,-1.9) {\small Class 2: F--L--E\\[-2pt]\small (standing on tiptoe)};
% Class 3
\draw[thick,popdark] (5.5,0) -- (9.5,0);
\fill[poporange] (9.3,-0.7) -- (8.9,0) -- (9.7,0) -- cycle;
\node[poporange] at (9.3,-1.1) {\small F};
\draw[->,popink,thick] (6.7,0.9) -- (6.7,0.1);
\node[popink] at (6.7,1.2) {\small E};
\draw[->,popgreen,thick] (5.7,0.9) -- (5.7,0.1);
\node[popgreen] at (5.7,1.2) {\small L};
\node[popdark,align=center] at (7.5,-1.9) {\small Class 3: L--E--F\\[-2pt]\small (flexing the elbow)};
\end{tikzpicture}
\legende{The three classes of levers. F = fulcrum, E = effort, L = load.}
\end{popfigure}

\begin{itemize}
\item \textbf{Class 1} (fulcrum between effort and load): nodding the head --- the skull pivots on the atlas (fulcrum), neck muscles provide the effort behind, and the weight of the face is the load in front.
\item \textbf{Class 2} (load between fulcrum and effort): rising on tiptoe --- the toes are the fulcrum, the body weight through the ankle is the load, and the calf muscles pulling the heel provide the effort. This class gives a \emph{mechanical advantage} (large load moved by smaller effort).
\item \textbf{Class 3} (effort between fulcrum and load): flexing the elbow --- the elbow is the fulcrum, the biceps pulls on the radius close to the joint, and the load is in the hand. This is the commonest arrangement in the body; it favours \emph{speed and range of movement} over force.
\end{itemize}

\courssub{Joint disorders and injuries}

Joints are put under great stress in daily life and sport, and several disorders are common.

\begin{center}
\begin{tabularx}{\linewidth}{|l|X|X|}
\hline
\rowcolor{popblueL} \textbf{Disorder} & \textbf{Cause} & \textbf{Prevention / care} \\
\hline
Sprain & Ligaments stretched or torn by sudden twisting of a joint (e.g. ankle on rough ground) & Warm up before sport; wear proper shoes; rest and ice the joint \\
\hline
Dislocation & Bone forced out of its normal position in the joint, often tearing ligaments and capsule & Avoid violent twists and falls; must be repositioned by trained personnel \\
\hline
Arthritis & Inflammation of a joint; cartilage may wear away (osteoarthritis) or the joint is attacked by the immune system (rheumatoid arthritis) & Healthy body weight, moderate exercise, medical treatment \\
\hline
\end{tabularx}
\end{center}

\begin{savaistu}[Joints and sport in Cameroon]
Ankle sprains and knee ligament tears are among the most frequent injuries of footballers playing on the hard, uneven pitches common in many Cameroonian towns and villages. Simple measures --- warming up, strengthening the muscles around joints, and playing on well-maintained grounds --- greatly reduce these injuries. This is why physiotherapy after a sprain is as important as the first aid itself.
\end{savaistu}

\begin{exoresolu}[Identifying joints and explaining smooth movement]
\textbf{Statement.} (a) Name the type of synovial joint found at the hip and at the elbow, and state one difference in the movements they allow. (b) Explain \emph{three} features of a synovial joint that make movement smooth and painless. (c) A student writes: ``Tendons hold the bones of the knee together.'' Correct this statement.
\tcblower
\textbf{Solution.} (a) The hip is a \textbf{ball-and-socket} joint, allowing movement in all planes including rotation; the elbow is a \textbf{hinge} joint, allowing movement in one plane only (flexion and extension).
(b) (i) \textbf{Articular cartilage} covers the bone ends, giving smooth surfaces that reduce friction and absorb shock. (ii) \textbf{Synovial fluid}, secreted by the synovial membrane, lubricates the surfaces and nourishes the cartilage. (iii) \textbf{Ligaments} hold the bones firmly together so the joint stays stable while still permitting movement (the capsule keeps the lubricating fluid in place).
(c) The statement is wrong: it is the \textbf{ligaments} that hold the bones of the knee together (bone to bone). Tendons join \emph{muscle to bone} and transmit the force of contraction.
\end{exoresolu}

\begin{exercice}[Joints and movement]
\begin{enumerate}
\item Define a joint and name the three classes of joints according to degree of movement, giving one example of each.
\item Draw a labelled diagram of a synovial joint and state the function of each labelled part.
\item Explain, using the biceps and triceps, what is meant by antagonistic muscles.
\item To which class of lever does the forearm belong when the elbow is flexed? Justify your answer by identifying fulcrum, effort and load.
\end{enumerate}
\end{exercice}

\begin{corrige}[Exercise 1]
\begin{enumerate}
\item A joint is a point where two or more bones meet. Fixed joints (skull sutures), slightly movable joints (between vertebrae, pubic symphysis), freely movable or synovial joints (shoulder, knee).
\item Diagram as in the figure above: bone ends covered by articular cartilage (reduces friction, absorbs shock); synovial membrane (secretes synovial fluid); synovial fluid in the cavity (lubricates and nourishes cartilage); capsule (encloses joint, retains fluid); ligaments (join bone to bone, stabilise the joint).
\item The biceps and triceps are antagonistic because they produce opposite movements at the elbow: contraction of the biceps (with relaxation of the triceps) flexes the arm; contraction of the triceps (with relaxation of the biceps) extends it.
\item Class 3 lever: the elbow joint is the fulcrum, the biceps insertion on the radius (between fulcrum and load) provides the effort, and the load is held in the hand; the effort lies between fulcrum and load.
\end{enumerate}
\end{corrige}

\begin{aretenir}[Key points]
\begin{itemize}
\item Joints are points where bones meet; they may be fixed, slightly movable, or freely movable (synovial).
\item Synovial joints are classified by the shape of their surfaces: hinge, ball-and-socket, pivot, gliding, saddle.
\item Smooth movement is ensured by articular cartilage, synovial fluid (lubrication and nutrition), the capsule and ligaments.
\item Ligaments join bone to bone; tendons join muscle to bone.
\item Muscles work in antagonistic pairs (flexor/extensor) because they can only pull.
\item Bones act as levers, with the joint as fulcrum, the muscle pull as effort and the resistance as load.
\end{itemize}
\end{aretenir}

\coursec{Muscle contraction}

In the previous section we saw that joints allow the bones of the skeleton to move smoothly past one another. But a joint cannot move by itself: something must supply the \cle{force}. That something is \textbf{muscle}. Every time you lift a hoe on the farm, kick a football in Yaoundé, or simply blink, thousands of microscopic protein filaments are sliding past one another inside your muscle cells, converting chemical energy from ATP into mechanical work. In this section we study the types of muscle tissue, the structure of a skeletal muscle from the whole organ down to the sarcomere, the sliding filament theory of contraction, its nervous and chemical control, and how the energy for contraction is supplied.

\courssub{Types of muscle tissue}

Observe three everyday situations: you decide to wave your hand (it obeys immediately); your heart beats whether you think about it or not; and food moves along your gut without any conscious command. These three situations correspond to the three types of muscle tissue found in the mammalian body.

\begin{definition}[Types of muscle tissue]
There are three types of \cle{muscle tissue}:
\begin{itemize}
\item \textbf{Skeletal (striated, voluntary) muscle}: attached to bones by tendons; under \cle{voluntary} control of the somatic nervous system; cells (fibres) are long, cylindrical, multinucleate and show light and dark \cle{striations} (bands) under the microscope.
\item \textbf{Cardiac muscle}: found only in the wall of the \cle{heart}; striated but \cle{involuntary}; cells are branched, usually uninucleate, joined end to end by intercalated discs, and contract rhythmically without fatigue.
\item \textbf{Smooth (unstriated, involuntary) muscle}: found in the walls of the gut, blood vessels, bladder and uterus; spindle-shaped uninucleate cells with no striations; contracts slowly and sustains tension for long periods.
\end{itemize}
\end{definition}

\begin{center}
\begin{tabularx}{\linewidth}{|l|X|X|X|}
\hline
\rowcolor{popblueL} \textbf{Feature} & \textbf{Skeletal} & \textbf{Cardiac} & \textbf{Smooth} \\
\hline
Striations & Present & Present & Absent \\
\hline
Control & Voluntary & Involuntary & Involuntary \\
\hline
Location & Attached to skeleton & Heart wall & Gut, vessels, uterus \\
\hline
Cell shape & Long cylindrical, multinucleate & Branched, 1--2 nuclei & Spindle-shaped, 1 nucleus \\
\hline
Speed of contraction & Fast, fatigues & Rhythmic, does not fatigue & Slow, sustained \\
\hline
\end{tabularx}
\end{center}

In this section we concentrate on \textbf{skeletal muscle}, because it is the effector of locomotion and the type examined at Advanced Level.

\courssub{Gross structure of a skeletal muscle}

If you dissect the muscle of a chicken leg (a common practical in our laboratories) and tease it apart, you discover a remarkable hierarchy of organisation, each level wrapped in connective tissue:

\begin{itemize}
\item The \textbf{whole muscle} (e.g.\ the biceps) has a thick central \cle{muscle belly} and is attached to bone at each end by a tough, inelastic \cle{tendon} made of collagen. The whole muscle is wrapped in a connective tissue sheath, the \emph{epimysium}.
\item Inside, the muscle is divided into bundles called \cle{fascicles}, each wrapped in the \emph{perimysium}.
\item Each fascicle contains many \cle{muscle fibres} — these are the actual muscle cells, up to several centimetres long, each wrapped in the \emph{endomysium}.
\item Inside each fibre lie hundreds of parallel \cle{myofibrils}, the thread-like structures that actually contract, bathed in \emph{sarcoplasm} (muscle cytoplasm) containing many mitochondria and a specialised smooth endoplasmic reticulum, the \cle{sarcoplasmic reticulum}, which stores calcium ions.
\item Each myofibril is a chain of repeating contractile units, the \cle{sarcomeres}.
\end{itemize}

\begin{popfigure}
\begin{tikzpicture}[scale=0.9, every node/.style={font=\small}]
% Whole muscle
\draw[fill=poporange!60, draw=popdark] (0,0.6) ellipse (0.45 and 0.9);
\draw[fill=poporange!60, draw=popdark] (0.45,0.6) ellipse (0.45 and 0.9);
\draw[fill=poporange!60, draw=popdark] (0.9,0.6) ellipse (0.45 and 0.9);
\draw[fill=gray!30] (-0.9,0.35) rectangle (-0.45,0.85);
\draw[fill=gray!30] (1.35,0.35) rectangle (1.8,0.85);
\node at (0.45,-0.7) {Whole muscle};
\node[font=\scriptsize] at (2.1,0.6) {tendon};
% Fascicle
\draw[fill=popgreen!40, draw=popdark] (4.2,0.6) ellipse (0.5 and 0.8);
\foreach \x/\y in {4.0/0.9,4.4/0.9,4.0/0.3,4.4/0.3,4.2/0.6}{\draw[fill=popgreen!70,draw=popdark] (\x,\y) circle (0.12);}
\node at (4.2,-0.7) {Fascicle};
% Fibre
\draw[fill=popblue!40, draw=popdark] (6.6,0.6) ellipse (0.35 and 0.7);
\draw[fill=popblue!70, draw=popdark] (6.6,0.6) ellipse (0.15 and 0.45);
\node at (6.6,-0.7) {Muscle fibre};
% Myofibril
\draw[fill=poppurple!30, draw=popdark] (8.6,0.2) rectangle (10.4,1.0);
\foreach \x in {8.9,9.5,10.1}{\draw[popdark,thick] (\x,0.2) -- (\x,1.0);}
\node at (9.5,-0.7) {Myofibril};
% arrows
\draw[->,thick,popdark] (1.9,0.6) -- (3.5,0.6);
\draw[->,thick,popdark] (4.9,0.6) -- (6.1,0.6);
\draw[->,thick,popdark] (7.1,0.6) -- (8.4,0.6);
\draw[->,thick,popdark] (9.5,-0.95) -- (9.5,-1.5);
% Sarcomere
\draw[popdark,very thick] (8.0,-2.3) -- (8.0,-1.6);
\draw[popdark,very thick] (11.0,-2.3) -- (11.0,-1.6);
\foreach \x in {8.35,8.55,8.75}{\draw[popblue,thick] (\x,-1.95) -- (9.3,-1.95);}
\foreach \x in {10.25,10.45,10.65}{\draw[popblue,thick] (9.7,-1.95) -- (\x,-1.95);}
\foreach \y in {-2.1,-1.95,-1.8}{\draw[poporange,very thick] (9.3,\y) -- (9.7,\y);}
\node[font=\scriptsize] at (8.0,-2.5) {Z};
\node[font=\scriptsize] at (11.0,-2.5) {Z};
\node at (9.5,-2.9) {Sarcomere};
\end{tikzpicture}
\legende{Levels of organisation of a skeletal muscle: whole muscle $\rightarrow$ fascicle $\rightarrow$ fibre $\rightarrow$ myofibril $\rightarrow$ sarcomere.}
\end{popfigure}

\courssub{Ultrastructure of the myofibril: the sarcomere}

Under the electron microscope, each myofibril is seen to consist of two kinds of protein filament:

\begin{itemize}
\item \textbf{Thick filaments}, made of the protein \cle{myosin}; each myosin molecule has a long tail and a globular \emph{head} that can bind to actin and hydrolyse ATP (it acts as an ATPase).
\item \textbf{Thin filaments}, made mainly of the protein \cle{actin}, together with the regulatory proteins \cle{troponin} and \cle{tropomyosin}.
\end{itemize}

These filaments are arranged in a precise repeating pattern that produces the striations:

\begin{itemize}
\item The \cle{Z-line} is a dark disc to which the thin filaments are anchored.
\item The \cle{I-band} is the light region containing thin filaments only.
\item The \cle{A-band} is the dark region spanning the whole length of the thick filaments (including the zone where thick and thin filaments overlap).
\item The \cle{H-zone} is the lighter centre of the A-band, containing thick filaments only.
\end{itemize}

\begin{definition}[The sarcomere]
A \cle{sarcomere} is the region of a myofibril between two successive \cle{Z-lines}. It is the \textbf{functional (contractile) unit} of a skeletal muscle fibre. A single myofibril contains thousands of sarcomeres arranged end to end, and because they all shorten together, their tiny individual shortenings add up to a large shortening of the whole muscle.
\end{definition}

\courssub{The sliding filament theory}

\begin{propriete}[The sliding filament theory]
During muscle contraction, the thin (actin) filaments \textbf{slide inwards} past the thick (myosin) filaments, pulling the Z-lines closer together. \textbf{Neither the actin nor the myosin filaments themselves shorten.} Consequently:
\begin{itemize}
\item the \cle{I-band} \textbf{shortens};
\item the \cle{H-zone} \textbf{shortens} (and may disappear);
\item the \cle{A-band} \textbf{stays the same length} (it equals the constant length of the thick filaments);
\item the \textbf{sarcomere shortens} and the Z-lines move closer together.
\end{itemize}
This theory (developed by H.\,E. Huxley, A.\,F. Huxley and co-workers from electron-microscope observations) is supported by electron micrographs of relaxed and contracted muscle, and you may be asked to interpret such micrographs in the examination.
\end{propriete}

\begin{popfigure}
\begin{tikzpicture}[scale=1.0, every node/.style={font=\small}]
% RELAXED
\node[font=\small\bfseries] at (3,2.6) {Relaxed sarcomere};
\draw[popdark,very thick] (0.5,0.2) -- (0.5,2.2);
\draw[popdark,very thick] (5.5,0.2) -- (5.5,2.2);
% thin filaments
\foreach \y in {0.5,0.9,1.3,1.7}{\draw[popblue,thick] (0.5,\y) -- (2.2,\y);}
\foreach \y in {0.5,0.9,1.3,1.7}{\draw[popblue,thick] (3.8,\y) -- (5.5,\y);}
% thick filaments
\foreach \y in {0.6,1.0,1.4,1.8}{\draw[poporange,very thick] (2.2,\y) -- (3.8,\y);}
% labels
\draw[<->,popmuted] (0.5,0.0) -- (2.2,0.0); \node[font=\scriptsize] at (1.35,-0.3) {I-band};
\draw[<->,popmuted] (2.2,0.0) -- (3.8,0.0); \node[font=\scriptsize] at (3.0,-0.3) {A-band};
\draw[<->,popmuted] (2.55,2.4) -- (3.45,2.4); \node[font=\scriptsize] at (3.0,2.65) {H-zone};
\node[font=\scriptsize] at (0.5,-0.6) {Z};
\node[font=\scriptsize] at (5.5,-0.6) {Z};
% CONTRACTED
\node[font=\small\bfseries] at (9.5,2.6) {Contracted sarcomere};
\draw[popdark,very thick] (7.2,0.2) -- (7.2,2.2);
\draw[popdark,very thick] (11.4,0.2) -- (11.4,2.2);
\foreach \y in {0.5,0.9,1.3,1.7}{\draw[popblue,thick] (7.2,\y) -- (8.9,\y);}
\foreach \y in {0.5,0.9,1.3,1.7}{\draw[popblue,thick] (9.7,\y) -- (11.4,\y);}
\foreach \y in {0.6,1.0,1.4,1.8}{\draw[poporange,very thick] (8.9,\y) -- (9.7,\y);}
\draw[<->,popmuted] (7.2,0.0) -- (8.9,0.0); \node[font=\scriptsize] at (8.05,-0.3) {I-band};
\draw[<->,popmuted] (8.9,0.0) -- (9.7,0.0); \node[font=\scriptsize] at (9.3,-0.3) {A-band};
\node[font=\scriptsize] at (7.2,-0.6) {Z};
\node[font=\scriptsize] at (11.4,-0.6) {Z};
% arrows showing sliding
\draw[->,popgreen,very thick] (8.0,2.35) -- (8.7,2.35);
\draw[->,popgreen,very thick] (10.6,2.35) -- (9.9,2.35);
\end{tikzpicture}
\legende{Relaxed and contracted sarcomeres. The thin filaments (blue) slide inwards over the thick filaments (orange): the I-band and H-zone shorten while the A-band is unchanged.}
\end{popfigure}

\begin{methode}[Explaining band changes as proof of sliding]
When an examination question shows electron micrographs of relaxed and contracted muscle and asks what they prove, proceed in numbered steps:
\begin{enumerate}
\item \textbf{Measure} (or state) that the A-band has the same length in both micrographs.
\item \textbf{Deduce} that the thick (myosin) filaments have not changed length, since the A-band equals the length of the thick filaments.
\item \textbf{Observe} that the I-band and H-zone are shorter (or absent) in the contracted muscle.
\item \textbf{Deduce} that the thin filaments have moved inwards, overlapping the thick filaments more, and that the Z-lines have been pulled closer together.
\item \textbf{Conclude}: the filaments \emph{slide} past each other; they do not themselves contract. This is the sliding filament theory.
\end{enumerate}
\end{methode}

\courssub{The cross-bridge cycle: how the filaments slide}

The sliding is produced by repeated cycles of attachment, pulling and detachment of the \cle{myosin heads}, called \textbf{cross-bridges}. Each cycle moves the thin filament by only a few nanometres, so the cycle is repeated many times per second during a single contraction.

\begin{methode}[The cross-bridge cycle]
\begin{enumerate}
\item \textbf{Attachment}: with the actin binding sites exposed (see control below), the energised myosin head binds to actin, forming a \cle{cross-bridge}.
\item \textbf{Power stroke}: the myosin head tilts, pulling the thin filament towards the centre of the sarcomere. ADP and inorganic phosphate (P$_i$) are released. This is the step that does mechanical work.
\item \textbf{Detachment}: a molecule of \cle{ATP} binds to the myosin head, causing it to \textbf{detach} from actin.
\item \textbf{Re-cocking}: the ATP is hydrolysed to ADP + P$_i$; the energy released returns the myosin head to its high-energy ("cocked") position, ready to attach again further along the actin filament.
\end{enumerate}
\end{methode}

\begin{popfigure}
\begin{tikzpicture}[scale=0.95, every node/.style={font=\small}]
% actin filament common
\foreach \i in {0,1,2,3}{
  \begin{scope}[shift={(0,-3.1*\i)}]
    \draw[popblue,very thick] (0,0) -- (5.2,0);
    \foreach \x in {0.4,1.0,1.6,2.2,2.8,3.4,4.0,4.6}{\draw[fill=popblue,draw=popblue] (\x,0) circle (0.09);}
  \end{scope}
}
% Stage 1: attach
\draw[poporange,very thick] (2.2,0.0) -- (2.6,0.7);
\draw[fill=poporange,draw=popdark] (2.7,0.85) circle (0.22);
\node[align=center,text width=3.2cm] at (6.9,0.4) {\textbf{1. Attach}\\ head binds actin (cross-bridge)};
% Stage 2: power stroke
\draw[poporange,very thick] (3.0,-3.1) -- (2.5,-2.4);
\draw[fill=poporange,draw=popdark] (2.42,-2.28) circle (0.22);
\draw[->,popgreen,very thick] (3.6,-3.5) -- (2.9,-3.5);
\node[align=center,text width=3.4cm] at (6.9,-2.7) {\textbf{2. Power stroke}\\ head tilts, pulls actin; ADP + P$_i$ released};
% Stage 3: detach
\draw[poporange,very thick] (2.6,-6.2) -- (2.6,-5.5);
\draw[fill=poporange,draw=popdark] (2.6,-5.35) circle (0.22);
\node[align=center,text width=3.4cm] at (6.9,-5.8) {\textbf{3. Detach}\\ ATP binds $\rightarrow$ head releases actin};
% Stage 4: recock
\draw[poporange,very thick] (1.8,-9.3) -- (2.3,-8.6);
\draw[fill=poporange,draw=popdark] (2.42,-8.45) circle (0.22);
\node[align=center,text width=3.4cm] at (6.9,-8.9) {\textbf{4. Re-cock}\\ ATP hydrolysed; head returns to high-energy position};
% cycle arrow
\draw[->,popdark,thick] (0.6,-9.9) .. controls (-0.8,-7) and (-0.8,-2) .. (0.6,-0.4);
\node[font=\scriptsize,align=center,text width=1.6cm] at (-1.0,-5.0) {cycle repeats};
\end{tikzpicture}
\legende{The cross-bridge cycle: attach, power stroke, detach (needs ATP), re-cock (ATP hydrolysed).}
\end{popfigure}

\begin{attention}[ATP is needed for detachment, not only for contraction]
A very common error is to write that ATP is used only "to make the muscle contract". In fact ATP has \textbf{two} roles: (i) its hydrolysis \emph{re-cocks} the myosin head, and (ii) the \emph{binding} of a fresh ATP molecule is what causes the myosin head to \textbf{detach} from actin. Without ATP, myosin heads cannot let go. This explains \cle{rigor mortis}: after death, respiration stops, ATP is exhausted, cross-bridges remain locked, and the body stiffens — until the muscle proteins later decompose. Rigor mortis is excellent evidence that ATP is required for \emph{relaxation} as well as contraction.
\end{attention}

\begin{attention}[Muscles can only pull, never push]
A muscle can exert force only by \textbf{shortening} (contracting); it cannot actively lengthen or push. This is why muscles work in \cle{antagonistic pairs}: the \textbf{biceps} contracts to flex the elbow, then the \textbf{triceps} contracts to extend it, stretching the relaxed biceps back to its original length. Never write "the biceps pushes the forearm down".
\end{attention}

\courssub{Control of contraction: nerve impulse to calcium release}

A muscle fibre contracts only when stimulated by a \textbf{motor neurone} at a specialised synapse, the \cle{neuromuscular junction}. The sequence is:

\begin{enumerate}
\item A \textbf{nerve impulse} (action potential) arrives at the motor neurone ending.
\item Vesicles release the transmitter \cle{acetylcholine}, which diffuses across the synaptic cleft and depolarises the muscle fibre membrane (\emph{sarcolemma}).
\item The action potential spreads over the sarcolemma and down inward folds called \cle{T-tubules}, deep into the fibre.
\item This triggers the \cle{sarcoplasmic reticulum} to release stored \textbf{calcium ions} (\ce{Ca^{2+}}) into the sarcoplasm.
\item \ce{Ca^{2+}} binds to \cle{troponin} on the thin filaments, causing \cle{tropomyosin} to shift position and \textbf{expose the myosin-binding sites} on actin.
\item Cross-bridge cycling begins and the sarcomere shortens.
\item When stimulation stops, \ce{Ca^{2+}} is actively pumped back into the sarcoplasmic reticulum (using ATP), tropomyosin covers the binding sites again, and the muscle \textbf{relaxes}.
\end{enumerate}

\begin{aretenir}[Key points of control]
\begin{itemize}
\item The \textbf{nerve impulse} triggers release of \ce{Ca^{2+}} from the sarcoplasmic reticulum.
\item \ce{Ca^{2+}} moves the \textbf{troponin--tropomyosin complex}, exposing the actin binding sites.
\item Relaxation is \emph{active}: pumping \ce{Ca^{2+}} back into the sarcoplasmic reticulum requires ATP.
\end{itemize}
\end{aretenir}

\courssub{Energy supply for contraction}

Muscle cannot store much ATP — only enough for a few seconds of activity. Three systems regenerate ATP, used in sequence as exercise continues:

\begin{itemize}
\item \textbf{Immediate (ATP + creatine phosphate)}: stored ATP is used first; then \cle{creatine phosphate} rapidly transfers its phosphate group to ADP:
\[ \text{creatine phosphate} + \text{ADP} \longrightarrow \text{creatine} + \text{ATP}. \]
This is very fast but is exhausted within roughly 10 seconds of maximal effort (e.g.\ a 100 m sprint).
\item \textbf{Anaerobic respiration (glycolysis)}: glucose is broken down to \cle{lactate} (lactic acid), yielding a small amount of ATP quickly without oxygen. It sustains intense effort for up to about 1--2 minutes, but lactate accumulates and contributes to \cle{fatigue}. Afterwards, extra oxygen is needed to oxidise the lactate and rebuild reserves — this is the \cle{oxygen debt}, which is why you keep breathing hard after a race.
\item \textbf{Aerobic respiration}: complete oxidation of glucose (and later fats) in the mitochondria; slow to start but produces \textbf{much more ATP per glucose molecule} and can continue for hours, as in the long-distance race from Buea to the summit of Mount Cameroon.
\end{itemize}

\begin{popfigure}
\begin{tikzpicture}
\begin{axis}[
  ybar, bar width=0.5cm,
  width=11cm, height=6.5cm,
  ymin=0, ymax=100,
  ylabel={Relative ATP supplied (\%)},
  symbolic x coords={0--10 s, 10 s--2 min, Over 2 min},
  xtick=data,
  x tick label style={font=\small},
  legend style={at={(0.5,1.02)}, anchor=south, legend columns=3, font=\scriptsize},
  nodes near coords, every node near coord/.append style={font=\scriptsize},
]
\addplot[fill=popgold, draw=popdark] coordinates {(0--10 s,80) (10 s--2 min,10) (Over 2 min,2)};
\addplot[fill=poporange, draw=popdark] coordinates {(0--10 s,15) (10 s--2 min,70) (Over 2 min,8)};
\addplot[fill=popgreen, draw=popdark] coordinates {(0--10 s,5) (10 s--2 min,20) (Over 2 min,90)};
\legend{Creatine phosphate, Anaerobic, Aerobic}
\end{axis}
\end{tikzpicture}
\legende{Approximate relative contribution of each energy pathway to ATP supply as exercise duration increases (typical values for maximal effort).}
\end{popfigure}

\courssub{Muscle fatigue and fibre types}

\textbf{Muscle fatigue} is the decline in force during prolonged activity, associated with depletion of ATP and glycogen, accumulation of lactate, and reduced \ce{Ca^{2+}} release from the sarcoplasmic reticulum. Not all fibres fatigue at the same rate:

\begin{center}
\begin{tabularx}{\linewidth}{|l|X|X|}
\hline
\rowcolor{popblueL} \textbf{Feature} & \textbf{Slow-twitch (red)} & \textbf{Fast-twitch (white)} \\
\hline
Contraction speed & Slow & Fast \\
\hline
Main respiration & Aerobic & Anaerobic (glycolytic) \\
\hline
Mitochondria / myoglobin & Many / much (red colour) & Few / little \\
\hline
Fatigue resistance & High (posture, endurance) & Low (sprinting, jumping) \\
\hline
\end{tabularx}
\end{center}

Regular training increases the number of mitochondria, the blood supply (capillarisation) and the glycogen stores in the fibres used — this is the physiological basis of \textbf{adaptation to exercise}.

\begin{exoresolu}[Interpreting sarcomere measurements]
\textbf{Statement.} In an electron micrograph of relaxed muscle, a sarcomere measures $2.4\ \mu\mathrm{m}$, the A-band $1.6\ \mu\mathrm{m}$ and the two I-band halves together contribute $0.8\ \mu\mathrm{m}$ to the sarcomere. In a contracted fibre the sarcomere measures $1.8\ \mu\mathrm{m}$. (a) What is the length of the A-band in the contracted fibre? (b) Calculate the new total length of the I-bands within one sarcomere. (c) Explain your answers.
\tcblower
\textbf{Solution.}
\begin{enumerate}
\item (a) The A-band equals the length of the thick filaments, which do not change length. Therefore the A-band is still $1.6\ \mu\mathrm{m}$.
\item (b) Sarcomere length $=$ A-band $+$ total I-band within it (the H-zone lies inside the A-band). So total I-band $= 1.8 - 1.6 = 0.2\ \mu\mathrm{m}$ (i.e.\ $0.1\ \mu\mathrm{m}$ at each end).
\item (c) The A-band is unchanged because the myosin filaments neither shorten nor lengthen. The I-band has shortened from $0.8\ \mu\mathrm{m}$ to $0.2\ \mu\mathrm{m}$ because the actin filaments have slid inwards over the myosin, increasing the zone of overlap (the H-zone has correspondingly shrunk). This is exactly what the sliding filament theory predicts.
\end{enumerate}
\end{exoresolu}

\begin{exercice}[Checking your understanding]
\begin{enumerate}
\item A patient is given a poison that blocks acetylcholine receptors at the neuromuscular junction. Explain, step by step, why the skeletal muscles cannot contract.
\item Explain why a corpse becomes stiff a few hours after death (rigor mortis), referring to myosin heads and ATP.
\item A sprinter and a marathon runner have different proportions of fibre types in their leg muscles. State and explain the differences you would expect.
\end{enumerate}
\end{exercice}

\begin{corrige}[Exercise 1]
\begin{enumerate}
\item The nerve impulse still arrives and acetylcholine is still released, but it cannot bind to its receptors, so the sarcolemma is not depolarised. No action potential travels along the T-tubules, so the sarcoplasmic reticulum does not release \ce{Ca^{2+}}. Without \ce{Ca^{2+}}, troponin does not move tropomyosin, the actin binding sites stay covered, no cross-bridges form, and the sarcomeres do not shorten.
\item After death, respiration stops and ATP is no longer produced. \ce{Ca^{2+}} leaks out of the sarcoplasmic reticulum, so cross-bridges form; but because detachment of myosin heads requires the binding of fresh ATP, the heads remain attached to actin. The muscles stay contracted and the body stiffens, until the muscle proteins eventually decompose.
\item The sprinter has a high proportion of fast-twitch fibres: they contract rapidly and powerfully using anaerobic glycolysis, ideal for a short burst, but they fatigue quickly. The marathon runner has mostly slow-twitch fibres, rich in mitochondria, myoglobin and capillaries, contracting aerobically and resisting fatigue over long periods.
\end{enumerate}
\end{corrige}

\begin{savaistu}[Muscles and malaria]
The violent shivering of a malaria attack — common in our hospitals during the rainy season — is rapid, involuntary contraction of skeletal muscles. Shivering is the body's way of generating heat through repeated cross-bridge cycling, which releases energy from ATP hydrolysis. It is a striking reminder that muscle contraction is not only for locomotion: it is also central to temperature regulation.
\end{savaistu}

\begin{aretenir}[Summary of muscle contraction]
\begin{itemize}
\item Skeletal muscle is organised hierarchically: muscle $\rightarrow$ fascicle $\rightarrow$ fibre $\rightarrow$ myofibril $\rightarrow$ \textbf{sarcomere} (between two Z-lines).
\item Contraction occurs by \textbf{sliding of actin past myosin}: the I-band and H-zone shorten, the A-band stays constant; the filaments themselves never shorten.
\item The \textbf{cross-bridge cycle} (attach, power stroke, detach, re-cock) is driven by ATP; ATP is also needed to \emph{detach} myosin — hence rigor mortis.
\item Nerve impulse $\rightarrow$ acetylcholine $\rightarrow$ \ce{Ca^{2+}} release $\rightarrow$ troponin--tropomyosin shifts $\rightarrow$ binding sites exposed $\rightarrow$ contraction.
\item ATP is regenerated by creatine phosphate (seconds), anaerobic respiration (minutes; lactate; oxygen debt) and aerobic respiration (hours).
\item Muscles can only \textbf{pull}; they work in antagonistic pairs.
\end{itemize}
\end{aretenir}

\coursec{Concepts of sustainable development}

In the previous sections of this chapter we studied how the skeleton, the joints and the muscles allow the human body to move, to work and to feed itself. But a body can only remain active if it is \emph{nourished, housed and kept healthy} — and all of that depends on the natural resources around us: fertile soils, forests, rivers, fish stocks and clean air. Just as a muscle cannot contract indefinitely without a continuous supply of glucose and oxygen, a nation cannot develop indefinitely by exhausting the resources that support it. This final section therefore widens our view from the individual organism to the whole of society, and asks a question that biology is uniquely placed to answer: \emph{how can humans meet their needs today without destroying the capacity of the Earth — and of Cameroon — to meet the needs of tomorrow?}

\courssub{Defining sustainable development}

\begin{definition}[Sustainable development (Brundtland, 1987)]
\cle{Sustainable development} is development that \textbf{meets the needs of the present without compromising the ability of future generations to meet their own needs}.
\end{definition}

This definition, proposed by the World Commission on Environment and Development (the ``Brundtland Commission'') in its 1987 report \emph{Our Common Future}, contains two essential ideas:
\begin{itemize}
\item \textbf{Needs}: development must first satisfy the essential needs of the world's poorest people — food, water, health, shelter, education.
\item \textbf{Limits}: the environment's capacity to meet present and future needs is limited by technology and by the state of ecosystems. Exceeding those limits today mortgages tomorrow.
\end{itemize}

The key intuition is that of \emph{living off the interest, not the capital}. A farmer who eats only the fruit of his trees can eat every year; a farmer who cuts down the trees to sell the wood eats well once — and then starves. Natural ecosystems are the \cle{capital}; the harvests, clean water and fertile soils they produce year after year are the \cle{interest}.

\courssub{The three pillars of sustainable development}

Sustainable development rests on \cle{three pillars} which must be satisfied \emph{simultaneously}:

\begin{itemize}
\item \textbf{Environmental protection}: ecosystems, biodiversity, climate and natural resources must be conserved so that they continue to function and to provide their services.
\item \textbf{Economic viability}: activities must create wealth and employment in a durable way; a project that destroys its own resource base is not viable.
\item \textbf{Social equity}: benefits and costs must be shared fairly, within the present generation (rich and poor, town and village) and between generations.
\end{itemize}

The three pillars are \textbf{interdependent}: none can stand alone. A nature reserve that impoverishes the surrounding villagers will be invaded and destroyed (environment fails without social equity); a profitable logging company that clears the whole forest collapses with the forest (economy fails without environment); an equitable but bankrupt society cannot fund schools, clinics or conservation (everything fails without a viable economy).

\begin{popfigure}
\begin{tikzpicture}[scale=1.0]
\fill[popgreenL,opacity=0.6] (90:1.1) circle (1.9);
\fill[popblueL,opacity=0.6] (210:1.1) circle (1.9);
\fill[poporange,opacity=0.25] (330:1.1) circle (1.9);
\draw[popgreen,thick] (90:1.1) circle (1.9);
\draw[popblue,thick] (210:1.1) circle (1.9);
\draw[poporange,thick] (330:1.1) circle (1.9);
\node[popgreen,font=\bfseries,align=center] at (90:2.6) {Environment};
\node[popblue,font=\bfseries,align=center] at (210:2.7) {Economy};
\node[poporange,font=\bfseries,align=center] at (330:2.7) {Society};
\node[popdark,font=\bfseries,align=center] at (0,0) {Sustain-\\ability};
\node[popmuted,align=center,font=\small] at (150:1.55) {bearable};
\node[popmuted,align=center,font=\small] at (30:1.55) {viable};
\node[popmuted,align=center,font=\small] at (270:1.7) {equitable};
\end{tikzpicture}
\legende{The three pillars of sustainable development. Only the central zone, where environment, economy and society overlap, is truly sustainable. The pairwise overlaps are described as bearable (environment + society), viable (environment + economy) and equitable (economy + society).}
\end{popfigure}

\begin{methode}[Analysing a development project through the three pillars]
To evaluate any project (a plantation, a dam, a fishing port, a new road) in an examination or in real life:
\begin{enumerate}
\item \textbf{Environment}: does the project degrade soil, water, forests or biodiversity? Can the resource regenerate? What waste is produced?
\item \textbf{Economy}: does it create durable income and jobs, or only short-term profit? Who receives the income?
\item \textbf{Society}: are local populations consulted and fairly compensated? Are the needs of women, children and future generations considered?
\item \textbf{Conclusion}: the project is sustainable \emph{only if all three answers are satisfactory}. If one pillar fails, propose modifications (replanting, quotas, benefit-sharing) rather than simply accepting or rejecting the project.
\end{enumerate}
\end{methode}

\begin{attention}[Sustainable development is not ``no development'']
A very common mistake is to confuse sustainable development with stopping all exploitation of nature. Sustainable development does \textbf{not} mean leaving forests untouched or refusing industrialisation; it means \emph{developing in a way that can continue indefinitely}. Cameroon needs roads, farms, factories and jobs — but they must be built so that soils, forests and fisheries are still productive for our grandchildren.
\end{attention}

\courssub{Renewable and non-renewable resources}

\begin{definition}[Renewable and non-renewable resources]
A \cle{non-renewable resource} (petroleum, minerals, fossil fuels) exists in a fixed stock that is exhausted by use. A \cle{renewable resource} (forests, fisheries, soils, fresh water, solar energy) is regenerated by natural processes and can, in principle, be used forever.
\end{definition}

The crucial biological point is that a biological resource is renewable \emph{only conditionally}:

\begin{propriete}[Condition of renewability]
A biological resource is renewable \textbf{only when the rate of harvesting does not exceed the rate of regeneration}. If exploitation exceeds regeneration, the stock declines and may collapse irreversibly (e.g. a fish population driven below its minimum viable size).
\end{propriete}

This property follows directly from population ecology. A population of fish or trees grows fastest at intermediate density (recall logistic growth from your ecology work): the \cle{maximum sustainable yield (MSY)} is obtained by harvesting only the surplus production, keeping the population near the size where its growth rate is maximal. Harvesting more than this surplus pushes the population down the curve; harvesting far too much drives it below the threshold from which it can recover, and on to extinction. Because regeneration varies from year to year, prudent managers apply a \cle{precautionary approach}: quotas are set \emph{below} the theoretical maximum to leave a safety margin.

\begin{popfigure}
\begin{tikzpicture}
\begin{axis}[
width=0.85\linewidth, height=6.5cm,
xlabel={Time (years)}, ylabel={Resource stock (tonnes)},
xmin=0, xmax=10, ymin=0, ymax=110,
xtick={0,2,4,6,8,10}, ytick={0,25,50,75,100},
legend style={at={(0.03,0.03)},anchor=south west,font=\small},
axis lines=left, thick]
\addplot[popgreen, very thick, domain=0:10, samples=60] {100 - 30*x/(1+0.35*x)};
\addlegendentry{Sustainable: harvest = regeneration}
\addplot[poporange, very thick, domain=0:10, samples=60] {100*exp(-0.28*x)};
\addlegendentry{Overexploitation: harvest $>$ regeneration}
\draw[popmuted, dashed] (axis cs:0,70) -- (axis cs:10,70);
\node[popmuted, font=\small, anchor=south] at (axis cs:5,70) {regeneration capacity};
\end{axis}
\end{tikzpicture}
\legende{Fate of a biological resource under two harvesting regimes. When the harvest equals the regeneration rate the stock stabilises (green); when the harvest exceeds regeneration the stock collapses (orange).}
\end{popfigure}

\courssub{Lessons from biology for sustainability}

Biology provides the scientific foundation of sustainable development:

\begin{itemize}
\item \textbf{Nutrient cycling}: in natural ecosystems, carbon, nitrogen and phosphorus circulate indefinitely; there is no ``waste''. Sustainable agriculture imitates this by recycling crop residues, manure and compost instead of depending solely on imported fertilisers.
\item \textbf{Carrying capacity}: every environment supports a maximum population of a given species. Human populations and livestock herds that exceed the carrying capacity of their land degrade it — overgrazing in the Far North of Cameroon is a clear example.
\item \textbf{Limits to population growth}: exponential growth cannot continue in a finite environment; planning (family, food production, urbanisation) must anticipate limits rather than collide with them.
\item \textbf{Ecosystem services}: forests regulate climate and rainfall, mangroves protect coasts and nurse fish, wetlands purify water, insects pollinate crops. These services are free but not infinite; destroying the ecosystem destroys the service.
\end{itemize}

\begin{savaistu}[Cameroon, a biological treasure]
Cameroon is often called ``Africa in miniature'' because it contains nearly all the continent's ecosystems, from Sahelian savanna to Congo-Basin rainforest. Its forests harbour exceptional biodiversity — gorillas, chimpanzees, forest elephants, and plant species found nowhere else. The \cle{Korup National Park} (South West Region) and the \cle{Dja Faunal Reserve} (inscribed on the UNESCO World Heritage List in 1987) protect some of Africa's oldest and richest rainforests. This biodiversity is both a global heritage and a potential source of durable income through ecotourism, research and non-timber forest products — \emph{if} it is managed sustainably.
\end{savaistu}

\courssub{Sustainable management of Cameroon's biological resources}

\begin{itemize}
\item \textbf{Community forests}: Cameroonian forestry law allows village communities to manage defined forest areas, harvesting timber and non-timber products under an approved management plan. Revenue stays in the village, giving people a direct interest in protecting the forest.
\item \textbf{Regulated fishing}: closed seasons, minimum mesh sizes and protection of nursery zones (estuaries, mangroves) along the coast from Limbe to Kribi allow fish stocks to reproduce before capture.
\item \textbf{Agroforestry}: combining trees (fruit trees, cocoa, oil palm) with food crops on the same plot maintains soil cover, recycles nutrients, diversifies income and reduces pressure to clear new forest.
\item \textbf{Protected areas}: national parks and reserves (Korup, Dja, Waza, Campo-Ma'an) conserve biodiversity, watersheds and genetic resources that may prove vital for medicine and agriculture.
\end{itemize}

\courssub{Threats to sustainability in Cameroon}

\begin{itemize}
\item \textbf{Deforestation}: clearing for farms, logging, fuelwood and roads destroys habitat, reduces rainfall infiltration and releases stored carbon.
\item \textbf{Soil erosion}: once the protective vegetation is removed, tropical rains wash away the thin fertile topsoil; yields fall, farmers clear new land, and the cycle accelerates.
\item \textbf{Bushmeat trade}: uncontrolled hunting empties forests of animals (``empty forest syndrome''), threatening both biodiversity and the protein supply of rural populations.
\item \textbf{Overfishing}: excessive and unselective fishing reduces catches, destroying the very resource on which fishing communities depend.
\item \textbf{Pollution}: untreated urban waste, plastics, pesticides and oil contaminate rivers and groundwater, harming health and aquatic life.
\end{itemize}

These threats form self-reinforcing \emph{degradation cycles}, as shown below.

\begin{popfigure}
\begin{tikzpicture}[node distance=1.5cm, every node/.style={draw=popblue, thick, rounded corners, fill=popblueL, align=center, text width=3.2cm, font=\small}, every path/.style={->, very thick, popdark}]
\node (def) {Deforestation (farming, logging, fuelwood)};
\node (ero) [right=of def] {Soil erosion by rain and wind};
\node (fer) [right=of ero] {Loss of soil fertility, falling yields};
\node (pov) [below=of fer] {Poverty and food insecurity};
\node (pre) [left=of pov] {Pressure to clear new land};
\draw (def) -- (ero);
\draw (ero) -- (fer);
\draw (fer) -- (pov);
\draw (pov) -- (pre);
\draw (pre) -- (def);
\node[draw=popgreen, fill=popgreenL, text width=4.2cm, below left=0.6cm and -1.2cm of ero] (brk) {Breaking the cycle: agroforestry, terracing, reforestation, alternative incomes};
\draw[popgreen, dashed] (brk.north) -- (ero.south);
\end{tikzpicture}
\legende{The degradation cycle linking deforestation to poverty. Sustainable practices (agroforestry, reforestation, alternative incomes) break the cycle at the erosion step.}
\end{popfigure}

\courssub{Sustainable development and health}

Health is both a goal and a condition of sustainable development: a population weakened by disease cannot work, learn or invest in the future. Health is therefore part of the \cle{human capital} of the nation. Its environmental pillars are:
\begin{itemize}
\item \textbf{Clean water and sanitation}: safe drinking water and proper latrines prevent cholera, typhoid and dysentery — diseases of polluted environments.
\item \textbf{Balanced nutrition}: diverse, sustainably produced food prevents malnutrition and deficiency diseases.
\item \textbf{Disease prevention}: protecting wetlands and managing waste reduces vectors; recall from your study of malaria that environmental management (draining stagnant water, sleeping under treated nets) complements medicine.
\end{itemize}

\courssub{Individual and collective actions}

\begin{itemize}
\item \textbf{Individuals}: apply the ``3 Rs'' — \cle{reduce} consumption, \cle{reuse} objects, \cle{recycle} materials; avoid waste of water, food and fuelwood; plant trees; refuse illegally hunted bushmeat.
\item \textbf{Farmers}: practise crop rotation, composting, agroforestry and contour ploughing to protect soils.
\item \textbf{Communities and the State}: environmental education in schools, community forests, enforcement of fishing and hunting regulations, waste management in cities.
\item \textbf{International frameworks}: in 2015 the United Nations adopted the \cle{Sustainable Development Goals (SDGs)} — 17 goals for 2030 covering poverty, health, education, clean water, climate action and life on land and below water. Cameroon, like all UN member states, has committed to them.
\end{itemize}

\begin{exoresolu}[Analysing a logging project]
\textbf{Statement.} A company proposes to exploit 5\,000 ha of rainforest near a village in the South Region, exporting all timber abroad. Using the three pillars, evaluate the project and propose improvements.
\tcblower
\textbf{Solution.}
\begin{enumerate}
\item \emph{Environment}: clear-felling 5\,000 ha destroys habitat, exposes soil to erosion and releases carbon. Unsustainable as proposed. Improvement: selective logging of mature trees only, with a management plan guaranteeing that the volume cut per year does not exceed natural regrowth, plus compulsory reforestation of degraded areas.
\item \emph{Economy}: exporting raw timber brings short-term profit but little local value. Improvement: local sawmill and carpentry workshop to create durable jobs and keep added value in Cameroon.
\item \emph{Society}: the village receives nothing and loses hunting and gathering grounds. Improvement: consult the community, share revenue (community forest model), maintain access to non-timber forest products.
\item \emph{Conclusion}: the project as proposed fails all three pillars; modified as above — selective harvest within regeneration limits, local processing, benefit-sharing — it can become sustainable.
\end{enumerate}
\end{exoresolu}

\begin{exercice}[Fishery and regeneration]
A coastal fish stock regenerates at a maximum rate of 800 tonnes per year. Three fleets propose annual catches of 600 t, 800 t and 1\,000 t. For each case, state whether the exploitation is sustainable and justify your answer. What long-term catch do you predict in the third case?
\end{exercice}

\begin{corrige}[Exercise 1]
600 t $<$ 800 t: sustainable, the stock can even increase slightly. 800 t $=$ regeneration: sustainable at the limit (maximum sustainable yield), but with no safety margin — a bad year (drought, disease) would unbalance it, so a precautionary quota below 800 t is wiser. 1\,000 t $>$ regeneration: unsustainable; the stock declines year after year, so the actual catch will fall below 1\,000 t and eventually collapse towards zero, destroying the fishery and the jobs it supports.
\end{corrige}

\begin{aretenir}[Key points]
\begin{itemize}
\item Sustainable development meets present needs without compromising the ability of future generations to meet theirs (Brundtland, 1987).
\item It rests on three interdependent pillars: environment, economy, society.
\item A biological resource is renewable only if the harvest rate does not exceed the regeneration rate; the maximum sustainable yield harvests only the surplus production, and a precautionary margin is wise.
\item Biology (nutrient cycles, carrying capacity, ecosystem services) provides the scientific basis of sustainability.
\item In Cameroon: community forests, regulated fishing, agroforestry and protected areas (Korup, Dja) are sustainable strategies; deforestation, bushmeat trade, erosion, overfishing and pollution are the main threats.
\item Health (clean water, sanitation, nutrition) is a pillar of human capital; individual actions (3 Rs) and collective frameworks (SDGs) complete the response.
\end{itemize}
\end{aretenir}

\begin{attention}[Exam tip]
In any GCE essay on sustainable development, \textbf{always illustrate with a local Cameroonian example} — a community forest, the coastal fishery, an agroforestry cocoa farm, or the Dja Reserve. A precise, correctly used local example distinguishes an excellent script from a vague one. And remember: sustainable development is development that can \emph{continue}, not development that is \emph{stopped}.
\end{attention}

\newpage

\coursec{Integration activity}

\courssub{Aims of the activity}

This integration activity closes the chapter by asking you to use, in a single coherent piece of work, everything you have learned about \cle{skeletons}, \cle{joints}, \cle{muscle contraction} and \cle{sustainable development}. By the end of the activity you should be able to:

\begin{itemize}
\item identify the main structures of a synovial joint (the knee) on a diagram and relate each structure to its function;
\item explain how smooth, low-friction movement is achieved at a joint and predict the consequences of damage to cartilage, ligaments or synovial fluid;
\item describe, step by step, how skeletal muscle contracts during sprinting, using the sliding filament theory, and discuss the energy sources that supply ATP;
\item apply the three pillars of sustainable development (environmental, economic, social) to a real resource-management decision;
\item use quantitative data (regeneration rate of a forest) to justify a harvesting plan;
\item argue, listen and defend a position in a plenary scientific debate, as expected of a GCE Advanced Level candidate.
\end{itemize}

\begin{attention}[How this activity will be assessed]
Your teacher will assess: (1) the accuracy of your biological explanations, (2) the correct use of data and calculations, (3) the quality of the link you make between the two parts of the dossier, and (4) your participation in the debate. Work in groups of four; every member must be able to defend every part of the group's answer.
\end{attention}

\courssub{Situation}

\textbf{Part A -- The injured footballer.}

Arrey, aged 19, is the best striker of a football academy in Buea, on the slopes of Mount Cameroon. During a regional championship match in Limbe, he sprinted for the ball, planted his right foot to change direction suddenly, and felt a sharp pain in his right knee, which swelled within an hour. At the district hospital, the doctor examined him and ordered an MRI scan. The report states:

\begin{quote}
\emph{``Rupture of the anterior cruciate ligament (ACL) of the right knee; tear of the medial meniscus; synovial fluid present in excess (effusion). Articular cartilage of the femoral condyles intact.''}
\end{quote}

The physiotherapist gives you a simplified diagram of the knee joint (reproduced below) and the following data collected from Arrey's training file:

\begin{center}
\begin{tabularx}{\linewidth}{|l|X|}
\hline
\rowcolor{popblueL} \textbf{Item} & \textbf{Data} \\
\hline
Position & Striker; repeated sprints of 10--40 m \\
\hline
Typical sprint duration & 3--6 s, maximal effort \\
\hline
Training load & 2 sessions/day, 6 days/week, little rest \\
\hline
Warm-up before matches & Often skipped when the bus arrives late \\
\hline
Knee structures injured & ACL (ligament), medial meniscus (cartilage) \\
\hline
Symptoms & Pain, swelling, knee ``gives way'', reduced range of movement \\
\hline
\end{tabularx}
\end{center}

\begin{popfigure}
\begin{tikzpicture}[scale=1.0]
% Femur
\draw[thick, fill=popblueL] (-1.6,3.2) -- (-1.6,1.2) arc (180:270:0.5) -- (-0.6,0.7) -- (0.6,0.7) arc (270:360:0.5) -- (1.6,1.2) -- (1.6,3.2) -- cycle;
\node at (0,2.6) {\textbf{Femur}};
% Cartilage on femur
\draw[thick, popgreen, fill=popgreenL] (-1.1,0.72) arc (180:270:0.5) -- (0.6,0.22) arc (270:360:0.5) -- cycle;
% Tibia
\draw[thick, fill=popblueL] (-1.3,-0.5) -- (-1.3,-2.6) -- (1.3,-2.6) -- (1.3,-0.5) -- cycle;
\node at (0,-1.9) {\textbf{Tibia}};
% Meniscus
\draw[thick, poporange, fill=popgold!40] (-1.0,-0.15) -- (-0.2,-0.15) -- (-0.2,-0.45) -- (-1.0,-0.45) -- cycle;
\draw[thick, poporange, fill=popgold!40] (0.2,-0.15) -- (1.0,-0.15) -- (1.0,-0.45) -- (0.2,-0.45) -- cycle;
% ACL
\draw[very thick, poppink] (-0.35,0.55) -- (0.15,-0.45);
% PCL
\draw[very thick, poppurple] (0.35,0.55) -- (-0.15,-0.45);
% Joint capsule
\draw[thick, dashed, popdark] (-1.7,1.0) -- (-1.7,-0.9) -- (1.7,-0.9) -- (1.7,1.0);
% Labels
\node[align=center, font=\small] at (-3.4,0.9) {Articular\\cartilage};
\draw[->, popgreen] (-2.2,0.75) -- (-0.9,0.5);
\node[align=center, font=\small] at (3.3,-0.3) {Meniscus};
\draw[->, poporange] (2.3,-0.3) -- (1.05,-0.3);
\node[align=center, font=\small] at (-3.4,-0.9) {ACL\\(ruptured)};
\draw[->, poppink] (-2.2,-0.75) -- (-0.15,0.0);
\node[align=center, font=\small] at (3.4,1.2) {Joint capsule\\(synovial fluid inside)};
\draw[->, popdark] (2.6,0.95) -- (1.7,0.4);
\node[align=center, font=\small] at (3.3,-1.6) {PCL};
\draw[->, poppurple] (2.5,-1.5) -- (0.15,-0.1);
\end{tikzpicture}
\legende{Simplified sagittal view of the right knee joint (ACL = anterior cruciate ligament, PCL = posterior cruciate ligament).}
\end{popfigure}

\textbf{Part B -- The village forest.}

The surgery and rehabilitation cost about 1\,500\,000 FCFA, which Arrey's family cannot afford. The village council of his home village, in the forest zone of the South-West Region, proposes to raise the money by selling timber from the 200-hectare community forest. An agricultural extension officer provides the following data:

\begin{center}
\begin{tabularx}{\linewidth}{|l|X|}
\hline
\rowcolor{popblueL} \textbf{Item} & \textbf{Data} \\
\hline
Area of community forest & 200 ha \\
\hline
Exploitable timber volume & 120 m$^3$ per ha (average) \\
\hline
Natural regeneration / regrowth & about 2 m$^3$ per ha per year \\
\hline
Market price of timber & 25\,000 FCFA per m$^3$ \\
\hline
Other uses of the forest & Firewood, bush mango (\emph{Irvingia}), medicinal plants, hunting, watershed protection for the village stream \\
\hline
\end{tabularx}
\end{center}

Some councillors want to fell 60 ha at once (``the boy needs the money now''); others suggest a slower plan. The council asks your class, as the best science students of the area, for advice.

\courssub{Tasks}

\textbf{Part A -- Support and locomotion}

\begin{enumerate}
\item On the knee diagram, name the type of joint and the type of skeleton it belongs to, and justify your answers.
\item Identify the two damaged structures and state the tissue type and normal function of each.
\item Using the symptoms described (swelling, knee ``giving way'', reduced movement), explain which function of each damaged structure has been lost.
\item Explain the role of \cle{synovial fluid} and \cle{articular cartilage} in smooth movement, and predict what may happen to the knee in the long term if the torn meniscus is not treated.
\item During a 4-second maximal sprint, which energy system supplies most of the ATP to Arrey's leg muscles, and why? Name the molecule that directly provides energy for contraction.
\item Describe, in numbered steps, the sliding filament mechanism of contraction of one sarcomere, from the arrival of the nerve impulse to the power stroke. Your answer must mention Ca$^{2+}$, troponin, tropomyosin, actin, myosin and ATP.
\item Explain why the ACL is particularly at risk when a player ``plants the foot and turns'', and why skipping the warm-up increases the risk of injury.
\item Propose a rehabilitation plan for Arrey (rest, immobilisation, physiotherapy, progressive return to training), justifying each stage biologically.
\end{enumerate}

\textbf{Part B -- Sustainable development}

\begin{enumerate}
\setcounter{enumi}{8}
\item Calculate the total exploitable timber volume of the forest and the money the village would obtain by felling all 200 ha.
\item Calculate the annual natural regrowth of the whole forest, and deduce the maximum volume that can be harvested \emph{each year} without reducing the forest capital. How much money per year does this represent?
\item Compare the two options (felling 60 ha at once vs.\ annual sustainable harvesting) in terms of the three pillars of sustainable development: environmental, economic and social.
\item The village needs 1\,500\,000 FCFA quickly. Propose a compromise harvesting plan that raises the money within two years while respecting the regeneration rate as far as possible. Show your calculations.
\item Give two additional measures (e.g.\ replanting, protecting the stream banks, controlling hunting) that would make the plan truly sustainable, and justify each one.
\item Prepare a 5-minute group presentation for the plenary debate: one member presents Part A, one presents Part B, one presents the compromise plan, and one answers questions from the ``village council'' (the rest of the class).
\end{enumerate}

\courssub{Model answer}

\textbf{Part A.}

\textbf{1--2.} The knee is a \cle{synovial joint} (a freely movable, hinge-type synovial joint) of the \cle{endoskeleton}: an internal bony skeleton characteristic of vertebrates, made of bone and cartilage, which grows with the body. The damaged structures are the anterior cruciate \cle{ligament} (dense fibrous connective tissue joining bone to bone, which stabilises the joint by limiting forward sliding and rotation of the tibia) and the medial \cle{meniscus} (fibrocartilage, which cushions the joint, spreads the load and improves the fit between femur and tibia).

\textbf{3.} The swelling is an effusion of synovial fluid caused by internal bleeding and inflammation. The knee ``giving way'' reflects the loss of ligament stability: without the ACL, the tibia slides abnormally under the femur. The reduced, painful range of movement reflects the torn meniscus interfering with the smooth gliding of the condyles.

\textbf{4.} Synovial fluid lubricates the joint and nourishes the cartilage (which has no blood supply); articular cartilage provides a smooth, low-friction, shock-absorbing surface. If the torn meniscus is not treated, load is concentrated on small areas of cartilage, friction and wear increase, and the cartilage of the femoral condyles -- still intact today -- will progressively degenerate (osteoarthritis), causing chronic pain and stiffness.

\textbf{5.} A 4-second maximal sprint is powered almost entirely by \cle{anaerobic} sources: the small store of ATP and \cle{phosphocreatine} (creatine phosphate) in the muscle fibres, which regenerates ATP very rapidly without oxygen. Aerobic respiration is too slow to reach full power in a few seconds, and anaerobic glycolysis to lactate contributes mainly in slightly longer efforts. The molecule that directly powers contraction is \cle{ATP}.

\textbf{6.} Sliding filament mechanism:
\begin{enumerate}
\item The nerve impulse arrives at the neuromuscular junction; acetylcholine triggers an action potential in the muscle fibre membrane, carried into the fibre by the T-tubules.
\item The sarcoplasmic reticulum releases Ca$^{2+}$ into the sarcoplasm.
\item Ca$^{2+}$ binds to \cle{troponin}, which moves \cle{tropomyosin} away from the myosin-binding sites on the \cle{actin} filaments.
\item Myosin heads (already energised by hydrolysis of ATP) bind to actin, forming cross-bridges.
\item The myosin heads pivot (power stroke), pulling the actin filaments towards the centre of the sarcomere; ADP and Pi are released.
\item A new ATP binds to each myosin head, detaching it from actin; ATP is hydrolysed and the cycle repeats.
\item The sarcomere shortens (H zone and I band shrink; A band constant); when Ca$^{2+}$ is pumped back into the sarcoplasmic reticulum, tropomyosin covers the binding sites again and the muscle relaxes.
\end{enumerate}

\textbf{7.} Planting the foot and turning fixes the tibia while the body rotates above it; the resulting twisting and forward-shearing force is exactly what the ACL resists, so it tears when the force exceeds its strength. Warming up raises muscle and connective-tissue temperature, increases blood flow, joint lubrication and tissue elasticity, and improves coordination; skipping it leaves stiff, poorly perfused tissues that tolerate sudden loads badly.

\textbf{8.} Rehabilitation plan: (i) acute phase -- rest, ice, compression and elevation to limit inflammation, because continued stress would enlarge the lesions; (ii) surgical repair or reconstruction of the ACL and meniscus, since ligaments and fibrocartilage heal poorly (poor blood supply); (iii) immobilisation followed by progressive mobilisation to restore synovial fluid circulation and prevent stiffness and cartilage degeneration; (iv) physiotherapy -- progressive strengthening of quadriceps and hamstrings, because strong muscles stabilise the joint and compensate for the ligament; (v) gradual return to training over several months, respecting tissue repair times, with systematic warm-up.

\textbf{Part B.}

\textbf{9.} Total volume $= 200 \times 120 = 24\,000\ \mathrm{m^3}$. Felling everything would bring $24\,000 \times 25\,000 = 600\,000\,000$ FCFA -- but only once, because the forest capital would be destroyed.

\textbf{10.} Annual regrowth $= 200 \times 2 = 400\ \mathrm{m^3}$ per year. Harvesting at most 400 m$^3$ per year leaves the forest capital intact; this yields $400 \times 25\,000 = 10\,000\,000$ FCFA per year, indefinitely.

\textbf{11.} Felling 60 ha at once: economically attractive in the short term ($60 \times 120 \times 25\,000 = 180\,000\,000$ FCFA), but environmentally destructive (loss of 7\,200 m$^3$ of standing capital, erosion on the slopes, loss of wildlife and of non-timber products) and socially damaging (the village loses firewood, food, medicines and stream protection; future generations inherit degraded land). Annual sustainable harvesting: environmentally sound, gives a permanent income, and preserves the social functions of the forest -- but cannot pay for the surgery immediately.

\textbf{12.} Compromise plan. The sustainable yield alone (400 m$^3$ per year, i.e.\ 10\,000\,000 FCFA) already covers the 1\,500\,000 FCFA needed, so only a small extra cut is required to build a safety margin. In year 1, harvest the full sustainable yield of 400 m$^3$ plus an extra 30 m$^3$ taken from a small, carefully chosen patch of about 0.25 ha: total $430\ \mathrm{m^3} \times 25\,000 = 10\,750\,000$ FCFA. In year 2, harvest only the regrowth (400 m$^3$, i.e.\ 10\,000\,000 FCFA). The two years raise $10\,750\,000 + 10\,000\,000 = 20\,750\,000$ FCFA, comfortably covering the surgery. The extra 30 m$^3$ represents less than 0.2\% of the forest capital ($30/24\,000 \approx 0.125\%$) and is compensated by replanting. The remaining money capitalises a village health fund fed by the annual 10\,000\,000 FCFA.

\textbf{13.} Additional measures: (i) replant indigenous tree species on any over-harvested patch to restore canopy and regeneration capacity; (ii) leave an unfelled buffer strip along the stream to prevent silting and protect the water supply; (iii) regulate hunting and the gathering of bush mango and medicinal plants so that non-timber resources are also renewed; (iv) involve the whole village, including women and youth, in managing the fund, so that the social pillar (equity, participation) is respected.

\textbf{Debate.} In the plenary, each group defends its plan; the ``council'' challenges the calculations and the biological reasoning. The expected conclusion: the same scientific rigour that explains how Arrey's muscles and joints must be allowed to \emph{regenerate} at their own pace also tells the village that a forest can only be harvested at the pace at which it \emph{regenerates} -- support, movement and sustainability all obey the logic of respecting living systems.

\begin{aretenir}[What this activity teaches]
The body and the forest are both renewable resources: ligaments, cartilage and muscle repair themselves \emph{if given time and the right conditions}, just as a forest regrows \emph{if harvesting never exceeds regeneration}. Sustainable development means meeting present needs -- Arrey's surgery -- without destroying the capital on which future needs depend.
\end{aretenir}

\newpage

\coursec{Methods and worked examples}

\courssub{Methods and worked examples}

In this part we transform the knowledge of the chapter into \cle{examination skills}. Each method card gives you a step-by-step reflex, and each worked example shows exactly how a GCE Advanced Level answer should be built and justified.

\begin{methode}[Identifying and comparing types of skeleton]
When a question asks you to identify a skeleton type or compare two skeletons:
\begin{enumerate}
\item \textbf{Locate the supporting structure}: is it a fluid-filled cavity (\cle{hydrostatic skeleton}, e.g. earthworm), a hard covering on the outside (\cle{exoskeleton}, e.g. crab, locust), or an internal framework of bone/cartilage (\cle{endoskeleton}, e.g. mammals, bony fish)?
\item \textbf{Check the material}: chitin (arthropods), calcium carbonate shells (molluscs), bone and cartilage (vertebrates).
\item \textbf{State how muscles attach}: muscles pull on the fluid wall, on the inside of the exoskeleton, or on the outside of bones.
\item \textbf{Compare growth}: an exoskeleton cannot grow, so arthropods must \cle{moult} (ecdysis); an endoskeleton grows with the body.
\item \textbf{Conclude with advantages and limitations} of each type (protection, weight, flexibility, size limit).
\end{enumerate}
\end{methode}

\begin{exoresolu}[Comparing the skeletons of an earthworm and a tilapia]
\textbf{Statement.} An earthworm (\emph{Lumbricus}) and a tilapia (a bony fish common in Cameroonian fish ponds) both move efficiently, yet they possess very different supporting structures. (a) Name the type of skeleton found in each organism. (b) Explain how each skeleton allows movement. (c) Give two advantages of the tilapia's skeleton over that of the earthworm.
\tcblower
\textbf{Solution.}
(a) The earthworm has a \textbf{hydrostatic skeleton}; the tilapia has a \textbf{bony endoskeleton}.

(b) \emph{Earthworm}: the body cavity (coelom) of each segment is filled with incompressible fluid. Circular and longitudinal muscles contract against this fluid, which acts as a firm but flexible ``skeleton''. When circular muscles contract, the segment becomes long and thin; when longitudinal muscles contract, it becomes short and fat. Alternating waves of contraction, anchored by bristles (chaetae), push the worm forward.

\emph{Tilapia}: an internal framework of bone (skull, vertebral column, fin rays) provides rigid levers. Blocks of muscle (myotomes) on either side of the vertebral column contract alternately, bending the flexible backbone from side to side and sweeping the tail fin through the water to produce thrust.

(c) Advantages of the endoskeleton:
\begin{itemize}
\item It \textbf{grows with the body}, so no moulting is needed and large body size is possible.
\item It provides \textbf{firm attachment points for muscles}, allowing more powerful, precise and rapid movements, and it protects delicate internal organs (e.g. the brain inside the skull).
\end{itemize}
\end{exoresolu}

\begin{methode}[Describing the functions of a named region of the vertebrate skeleton]
For questions such as ``Describe the structure and functions of the vertebral column'':
\begin{enumerate}
\item \textbf{Define the region} and state the number of bones (e.g. in humans: 7 cervical, 12 thoracic, 5 lumbar, 5 fused sacral, 4 fused coccygeal vertebrae).
\item \textbf{Relate structure to function}: e.g. the atlas supports the skull and allows nodding; the axis allows rotation; thoracic vertebrae articulate with ribs; lumbar vertebrae are massive to bear weight.
\item \textbf{List the general functions}: support, protection (spinal cord, brain, heart and lungs), muscle attachment, movement (leverage), blood cell production in bone marrow, and mineral storage (calcium, phosphorus).
\item Always \textbf{link each function to a structural feature} --- examiners award marks for the link, not for lists.
\end{enumerate}
\end{methode}

\begin{exoresolu}[Structure--function relationships in the human skeleton]
\textbf{Statement.} (a) State four functions of the mammalian skeleton. (b) For each function, name one part of the human skeleton adapted to it and explain the adaptation.
\tcblower
\textbf{Solution.}
\begin{center}
\begin{tabularx}{\linewidth}{|l|X|X|}
\hline
\rowcolor{popblueL} \textbf{Function} & \textbf{Part} & \textbf{Adaptation} \\
\hline
Support & Vertebral column, leg bones & Lumbar vertebrae are large and massive to bear the weight of the upper body; long bones have a rigid shaft of compact bone. \\
\hline
Protection & Skull, rib cage, vertebral column & The cranium is a fused bony box enclosing the brain; ribs form a cage around the heart and lungs; the neural arches enclose the spinal cord. \\
\hline
Movement & Limbs and girdles & Long bones act as levers; the pectoral and pelvic girdles give firm attachment to limb muscles; joints allow movement in defined directions. \\
\hline
Blood cell formation & Bone marrow of long bones, ribs, sternum & Red bone marrow in the spongy bone at the ends of long bones contains stem cells that produce red and white blood cells. \\
\hline
\end{tabularx}
\end{center}
(Any four of these rows earn full marks; note that each answer links the function to a named structure and its adaptation.)
\end{exoresolu}

\begin{methode}[Analysing a joint: type, movement and structures for smooth movement]
\begin{enumerate}
\item \textbf{Classify the joint}: fibrous/immovable (skull sutures), cartilaginous/slightly movable (between vertebrae), or \cle{synovial}/freely movable.
\item \textbf{If synovial, name the subtype and the movement}: hinge (elbow, knee --- one plane), ball-and-socket (hip, shoulder --- all planes), pivot (atlas--axis --- rotation), gliding (wrist).
\item \textbf{Draw or describe the structures ensuring smooth movement}: smooth \cle{articular cartilage} covering bone ends; \cle{synovial fluid} as lubricant; synovial membrane secreting the fluid; joint capsule holding everything together; ligaments binding bone to bone; tendons attaching muscle to bone.
\item \textbf{Explain each structure's role} in reducing friction or giving stability.
\end{enumerate}
\end{methode}

\begin{exoresolu}[The knee joint]
\textbf{Statement.} The diagram of a knee joint shows: the femur and tibia with cartilage at their ends, a capsule lined by a synovial membrane, synovial fluid, ligaments, and the tendon of the quadriceps muscle. (a) State the type of joint and the movement it allows. (b) Explain how three named structures ensure smooth, friction-free movement. (c) A footballer in Yaoundé tears a ligament. Explain why the joint becomes unstable.
\tcblower
\textbf{Solution.}
(a) The knee is a \textbf{synovial hinge joint}; it allows movement in one plane only --- flexion (bending) and extension (straightening) of the leg.

(b)
\begin{itemize}
\item \textbf{Articular cartilage}: a smooth, slippery layer covering the ends of the femur and tibia; it prevents the bones from rubbing directly on each other and absorbs shock.
\item \textbf{Synovial fluid}: an oily fluid secreted by the synovial membrane; it lubricates the cartilage surfaces, reducing friction to a minimum, and nourishes the cartilage.
\item \textbf{Joint capsule and ligaments}: the capsule encloses the fluid so it stays around the joint, while ligaments hold the bones in correct alignment so that movement occurs only in the intended plane.
\end{itemize}

(c) Ligaments bind bone to bone and prevent excessive or sideways movement. A torn ligament can no longer hold the femur and tibia firmly together, so the bones can slip out of alignment; the joint becomes loose and unstable, and may dislocate easily. This shows the essential stabilising role of ligaments.
\end{exoresolu}

\begin{methode}[Explaining how antagonistic muscles move a limb]
\begin{enumerate}
\item Identify the \textbf{pair of antagonistic muscles} (e.g. biceps = flexor, triceps = extensor at the elbow).
\item State that muscles can only \textbf{pull, never push}; bones act as levers and joints as fulcrums.
\item Describe the movement: when the flexor \textbf{contracts} and the extensor \textbf{relaxes}, the limb bends; the reverse straightens it.
\item Mention the role of \textbf{tendons} transmitting the force to the bone, and the nervous coordination ensuring one muscle relaxes while the other contracts (reciprocal inhibition).
\end{enumerate}
\end{methode}

\begin{exoresolu}[Lifting a load of plantains]
\textbf{Statement.} A market trader in Douala lifts a basket of plantains by bending her arm at the elbow. Describe the events in the muscles, bones and joint of the arm that bring about this movement.
\tcblower
\textbf{Solution.}
\begin{enumerate}
\item Nerve impulses from the motor neurones stimulate the \textbf{biceps (flexor)} to contract; simultaneously the \textbf{triceps (extensor)} is inhibited and relaxes.
\item The biceps pulls through its \textbf{tendon} on the radius (forearm bone).
\item The forearm acts as a \textbf{lever}, pivoting about the \textbf{elbow joint (fulcrum)}, a hinge joint.
\item The forearm is raised --- flexion occurs --- and the basket is lifted.
\item To lower the basket, the triceps contracts and the biceps relaxes, extending the arm. This opposite action is why the two muscles are called an \textbf{antagonistic pair}.
\end{enumerate}
Key idea: muscles only exert force by contracting (pulling); a second muscle is always needed to reverse the movement.
\end{exoresolu}

\begin{methode}[Interpreting the sliding filament theory from data or diagrams]
\begin{enumerate}
\item \textbf{Identify the bands} on a sarcomere diagram: A band (dark, thick myosin filaments), I band (light, actin only), H zone (myosin only), Z lines (boundaries of the sarcomere).
\item \textbf{Compare relaxed vs contracted}: during contraction the \textbf{H zone and I bands shorten}, the \textbf{A band stays the same length}, and the Z lines move closer together --- proof that filaments \emph{slide} rather than shorten.
\item \textbf{Sequence the mechanism}: nerve impulse $\rightarrow$ release of \ce{Ca^{2+}} from sarcoplasmic reticulum $\rightarrow$ \ce{Ca^{2+}} binds troponin, tropomyosin shifts, exposing actin binding sites $\rightarrow$ myosin heads attach (cross-bridges) $\rightarrow$ power stroke pulls actin towards the centre (ATP provides energy) $\rightarrow$ ATP detaches and re-cocks the myosin head $\rightarrow$ cycle repeats.
\item Remember: \textbf{ATP is needed both to detach myosin heads and to pump \ce{Ca^{2+}} back}; without ATP, muscles stay contracted (rigor mortis).
\end{enumerate}
\end{methode}

\begin{exoresolu}[Analysing sarcomere measurements]
\textbf{Statement.} A student measures the bands of a muscle sarcomere before and during contraction:

\begin{center}
\begin{tabular}{|l|c|c|}
\hline
\rowcolor{popblueL} \textbf{Band} & \textbf{Relaxed} & \textbf{Contracted} \\
\hline
A band & $1.6\ \mu\mathrm{m}$ & $1.6\ \mu\mathrm{m}$ \\
I band & $1.0\ \mu\mathrm{m}$ & $0.4\ \mu\mathrm{m}$ \\
H zone & $0.5\ \mu\mathrm{m}$ & $0.1\ \mu\mathrm{m}$ \\
\hline
\end{tabular}
\end{center}

(a) Explain how these results support the sliding filament theory. (b) Describe the role of calcium ions and ATP in bringing about the contraction.
\tcblower
\textbf{Solution.}
(a) The \textbf{A band is unchanged} ($1.6\ \mu\mathrm{m}$), showing that the thick myosin filaments do not shorten. The \textbf{I band and H zone both decrease}, showing that the thin actin filaments have \emph{slid inwards} between the myosin filaments, overlapping them more. Since the filaments themselves keep their length while the sarcomere shortens, contraction must occur by filaments \textbf{sliding past each other} --- exactly what the sliding filament theory states.

(b) \emph{Calcium ions}: when a nerve impulse reaches the muscle fibre, \ce{Ca^{2+}} is released from the sarcoplasmic reticulum. \ce{Ca^{2+}} binds to \textbf{troponin} on the actin filaments, causing \textbf{tropomyosin} to move aside and expose the myosin-binding sites on actin.

\emph{ATP}: myosin heads bind to the exposed sites forming cross-bridges; the heads then tilt (power stroke), pulling actin towards the centre of the sarcomere. ATP then binds to the myosin head, causing it to \textbf{detach}, and ATP hydrolysis re-cocks the head so the cycle can repeat. ATP also powers the pumps that return \ce{Ca^{2+}} to the sarcoplasmic reticulum for relaxation.
\end{exoresolu}

\begin{methode}[Distinguishing muscle types and energy supply for contraction]
\begin{enumerate}
\item \textbf{Classify the muscle}: skeletal (striated, voluntary, attached to bones), smooth (unstriated, involuntary, walls of gut and blood vessels), cardiac (striated, involuntary, branched, heart only, never fatigues).
\item \textbf{Energy supply}: immediate energy from ATP; regenerated by \textbf{aerobic respiration} (glucose + oxygen, in mitochondria) during moderate activity, and by \textbf{anaerobic respiration} (lactic acid produced, oxygen debt) during intense effort; phosphocreatine provides a rapid short-term ATP store.
\item \textbf{Link fibre type to activity}: slow-twitch (red, many mitochondria, endurance) vs fast-twitch (white, powerful bursts).
\end{enumerate}
\end{methode}

\begin{exoresolu}[Why a sprinter pants after a 100 m race]
\textbf{Statement.} After running a 100 m sprint, an athlete breathes deeply and rapidly for several minutes. (a) Explain the source of energy for the sprint. (b) Explain why the deep breathing continues after the race has ended.
\tcblower
\textbf{Solution.}
(a) During a maximal sprint, oxygen cannot be delivered to the leg muscles fast enough for aerobic respiration alone. ATP is supplied by (i) the small store of ATP and \textbf{phosphocreatine} in the muscle, and (ii) \textbf{anaerobic respiration} of glucose, which yields ATP rapidly but produces \textbf{lactic acid} as a waste product.

(b) The lactic acid accumulated in the muscles is carried by the blood to the liver, where it must be oxidised or converted back to glucose --- processes that \textbf{require oxygen}. In addition, ATP and phosphocreatine stores must be rebuilt by aerobic respiration. The extra oxygen needed after the exercise is called the \textbf{oxygen debt}. Continued deep, rapid breathing supplies this oxygen and removes the extra carbon dioxide produced, until the lactic acid is cleared and the stores are restored; breathing then returns to normal.
\end{exoresolu}

\begin{methode}[Applying concepts of sustainable development to a local problem]
\begin{enumerate}
\item \textbf{Define sustainable development}: development that meets the needs of the present generation \textbf{without compromising the ability of future generations} to meet their own needs.
\item Identify the \textbf{three pillars}: economic viability, social equity, and \textbf{environmental protection}.
\item Analyse the problem (deforestation, overfishing, bush burning, plastic pollution, overgrazing) in terms of resource depletion and ecosystem damage.
\item Propose \textbf{concrete sustainable solutions}: reforestation and agroforestry, fishing quotas and closed seasons, terracing and crop rotation, renewable energy, waste recycling, environmental education, protected areas (national parks).
\item Conclude by showing how the solution balances the three pillars.
\end{enumerate}
\end{methode}

\begin{exoresolu}[Deforestation around Mount Cameroon]
\textbf{Statement.} Communities on the slopes of Mount Cameroon clear forest each year for farms and firewood, causing loss of biodiversity and soil erosion. Using the concept of sustainable development, (a) explain why this practice is unsustainable, and (b) propose three measures that would allow the communities to meet their needs while protecting the forest.
\tcblower
\textbf{Solution.}
(a) The practice is unsustainable because the forest is being consumed \textbf{faster than it can regenerate}. Soil erosion removes fertile topsoil, so future harvests decline; loss of tree cover destroys habitats, reducing the unique \textbf{biodiversity} of Mount Cameroon; and the loss of vegetation reduces rainfall infiltration and increases flooding. The needs of the present generation (farmland, firewood) are being met by destroying the resources that \textbf{future generations} will depend on --- the exact opposite of sustainable development.

(b) Sustainable measures:
\begin{itemize}
\item \textbf{Agroforestry and reforestation}: planting trees alongside crops (e.g. fruit trees, nitrogen-fixing species) restores soil fertility, provides firewood on farms, and reduces the need to clear new forest --- protecting the \emph{environmental} pillar while keeping farms productive (\emph{economic} pillar).
\item \textbf{Improved cooking stoves and alternative fuels}: fuel-efficient stoves or biogas reduce firewood consumption, lowering pressure on the forest while saving families money and time (\emph{social} pillar).
\item \textbf{Community forest management and protected zones}: involving villagers in managing buffer zones around the Mount Cameroon National Park, with controlled harvesting and environmental education, ensures that resources are used at a rate that allows natural renewal, and that benefits are shared fairly within the community.
\end{itemize}
Each measure balances economic needs, social well-being and environmental protection --- the three pillars of sustainable development.
\end{exoresolu}

\begin{attention}[Common mistakes in this chapter]
\begin{itemize}
\item Saying that filaments ``shorten'' during contraction --- they \emph{slide}; only the sarcomere shortens.
\item Confusing \textbf{ligaments} (bone to bone) with \textbf{tendons} (muscle to bone).
\item Claiming the A band changes length --- it stays constant; the I band and H zone shrink.
\item Writing that muscles ``push'' --- muscles can only pull by contracting; limbs move through antagonistic pairs.
\item Confusing respiration with breathing, or saying lactic acid is produced because the muscle ``has no glucose'' --- it is produced when oxygen supply is insufficient.
\item Defining sustainable development vaguely as ``protecting nature'' --- you must mention \textbf{present and future generations} and the balance of economic, social and environmental needs.
\end{itemize}
\end{attention}

\newpage

\coursec{Exercises}

\courssub{I check my knowledge}

\begin{exercice}[Multiple choice questions]
For each question, choose the single correct answer.
\begin{enumerate}
\item The type of skeleton found in an earthworm is:
\begin{enumerate}
\item an exoskeleton;
\item an endoskeleton;
\item a hydrostatic skeleton;
\item a cartilaginous skeleton.
\end{enumerate}
\item The joint between the humerus and the scapula is:
\begin{enumerate}
\item a hinge joint;
\item a ball-and-socket joint;
\item a fixed joint;
\item a gliding joint.
\end{enumerate}
\item During muscle contraction, the region of the sarcomere that does \textbf{not} change in length is:
\begin{enumerate}
\item the I-band;
\item the H-zone;
\item the A-band;
\item the Z-line interval.
\end{enumerate}
\item The immediate source of energy for the detachment of myosin heads from actin is:
\begin{enumerate}
\item glucose;
\item lactic acid;
\item ATP;
\item creatine phosphate.
\end{enumerate}
\end{enumerate}
\end{exercice}

\begin{exercice}[True or false — justify]
State whether each of the following statements is true or false, and justify your answer in one or two sentences.
\begin{enumerate}
\item Cartilage at the ends of bones in a synovial joint reduces friction and absorbs shock.
\item Muscles produce movement by actively pushing against bones.
\item The biceps and triceps of the upper arm form an antagonistic pair.
\item Sustainable development means stopping all exploitation of natural resources.
\end{enumerate}
\end{exercice}

\begin{exercice}[Key definitions]
Define each of the following terms precisely, as you would in a GCE answer:
\begin{enumerate}
\item synovial joint;
\item sarcomere;
\item oxygen debt;
\item sustainable development.
\end{enumerate}
\end{exercice}

\begin{exercice}[Direct calculations]
\begin{enumerate}
\item A sarcomere is $2.2\ \mu\mathrm{m}$ long when relaxed and $1.6\ \mu\mathrm{m}$ long when fully contracted. Calculate the percentage decrease in length.
\item A muscle fibre contains $1\,500$ sarcomeres arranged end to end. Using the values in (a), calculate the total shortening of the fibre during maximal contraction, giving your answer in millimetres.
\item During strenuous exercise, an athlete's blood lactic acid rises from $1.0\ \mathrm{mmol\,dm^{-3}}$ at rest to $8.5\ \mathrm{mmol\,dm^{-3}}$. Calculate the increase and express it as a multiple of the resting value.
\end{enumerate}
\end{exercice}

\courssub{I practise}

\begin{exercice}[Labelling the human skeleton]
The diagram below shows the anterior view of the human skeleton with leader lines A--H pointing to: the cranium, the vertebral column, the rib cage, the humerus, the radius and ulna, the pelvic girdle, the femur, and the tibia and fibula.
\begin{enumerate}
\item Copy the diagram and label parts A to H.
\item State \textbf{one} function of each labelled region (support, protection, movement, or blood cell production).
\item Name the type of joint found between the femur and the pelvic girdle.
\end{enumerate}
\end{exercice}

\begin{exercice}[Comparing joints]
\begin{enumerate}
\item Construct a table comparing a hinge joint and a ball-and-socket joint under the following headings: example in the human body; planes of movement permitted; types of movement possible.
\item Explain why the shoulder joint, although very mobile, is more easily dislocated than the elbow joint.
\item Name the structures that hold the bones together at a synovial joint and state their role.
\end{enumerate}
\end{exercice}

\begin{exercice}[Smooth movement at a synovial joint]
\begin{enumerate}
\item Draw a large, fully labelled diagram of a typical synovial joint (e.g. the knee), showing: articular cartilage, synovial membrane, synovial fluid, joint capsule and ligaments.
\item For each labelled structure, state precisely how it contributes to smooth, friction-free movement.
\item Explain why a person whose synovial membrane produces too little synovial fluid experiences pain and stiffness during movement.
\end{enumerate}
\end{exercice}

\begin{exercice}[Interpreting lactic acid data]
The table below shows the blood lactic acid concentration of a student before, during and after a 400 m sprint.
\begin{center}
\begin{tabularx}{\linewidth}{|l|X|X|X|X|X|X|X|}
\hline
\rowcolor{popblueL} Time (min) & 0 & 2 & 4 & 6 & 10 & 20 & 30 \\
\hline
Lactic acid ($\mathrm{mmol\,dm^{-3}}$) & 1.0 & 3.5 & 7.8 & 8.2 & 6.0 & 2.5 & 1.1 \\
\hline
\end{tabularx}
\end{center}
(The sprint lasted from minute 0 to minute 4.)
\begin{enumerate}
\item Plot a graph of lactic acid concentration against time.
\item Explain the rise in lactic acid concentration between minutes 0 and 4.
\item Explain why lactic acid concentration remains high for some minutes \emph{after} the sprint has ended.
\item Use the graph to explain the term \cle{oxygen debt} and state why the student continues to breathe deeply after stopping.
\end{enumerate}
\end{exercice}

\courssub{I prepare for the GCE}

\begin{exercice}[Structured question: support and movement (GCE Paper 2 style) — 20 marks]
\begin{enumerate}
\item State three functions of the mammalian skeleton. \hfill (3 marks)
\item Name the bones that form the forelimb of a mammal, from the pectoral girdle to the digits. \hfill (4 marks)
\item Describe the structure of a synovial joint and explain how friction is minimised during movement. \hfill (6 marks)
\item Using the sliding filament theory, explain why, during contraction of a myofibril, the I-band and H-zone shorten while the A-band remains constant in length. \hfill (5 marks)
\item State the role of calcium ions and ATP in muscle contraction. \hfill (2 marks)
\end{enumerate}
\end{exercice}

\begin{exercice}[Essay: muscles, antagonism and levers — 20 marks]
``Muscles can only pull, never push.''

Discuss this statement with reference to:
\begin{itemize}
\item the mechanism of muscle contraction;
\item the arrangement of antagonistic muscle pairs, using the human arm as an example;
\item the action of bones as levers operated at joints.
\end{itemize}
Your essay should include at least one clearly labelled diagram. \hfill (20 marks)
\end{exercice}

\begin{exercice}[Essay: sustainable development in Cameroon — 20 marks]
\begin{enumerate}
\item Define the term \emph{sustainable development} and name its three pillars. \hfill (4 marks)
\item Using \textbf{one named Cameroonian natural resource} (for example the forests of the South Region, the fisheries of the coast, or the soils of the Western Highlands), explain how it can be exploited in a way that meets present needs without compromising the needs of future generations. \hfill (8 marks)
\item Discuss the interdependence of the environmental, economic and social pillars of sustainable development in a rural Cameroonian community, illustrating your answer with concrete examples. \hfill (8 marks)
\end{enumerate}
\end{exercice}

\newpage

\coursec{Solutions to the exercises}

\begin{corrige}[Exercise 1 — Multiple choice questions]
\begin{enumerate}
\item \textbf{Answer: (c) a hydrostatic skeleton.} The earthworm has no hard parts; its body wall muscles (circular and longitudinal) act against the incompressible coelomic fluid, which transmits pressure and gives the body its shape and rigidity.
\item \textbf{Answer: (b) a ball-and-socket joint.} The rounded head of the humerus fits into the shallow glenoid cavity of the scapula, permitting movement in all planes (flexion, extension, abduction, adduction, rotation).
\item \textbf{Answer: (c) the A-band.} The A-band corresponds to the length of the thick (myosin) filaments, which do not themselves shorten; only the I-band and H-zone narrow as the thin filaments slide inwards.
\item \textbf{Answer: (c) ATP.} The binding of a fresh ATP molecule to the myosin head causes it to detach from actin; hydrolysis of this ATP then re-cocks the head for the next cycle.
\end{enumerate}
\end{corrige}

\begin{corrige}[Exercise 2 — True or false — justify]
\begin{enumerate}
\item \textbf{True.} Articular (hyaline) cartilage covering the bone ends is smooth and slightly compressible: it provides a low-friction gliding surface and cushions the joint against mechanical shocks.
\item \textbf{False.} Muscle fibres can only exert force by \emph{shortening} (pulling on their points of attachment via tendons); they cannot actively lengthen or push. Movement in the opposite direction requires an antagonist muscle or gravity.
\item \textbf{True.} When the biceps contracts it flexes the elbow while the triceps relaxes; when the triceps contracts it extends the elbow while the biceps relaxes. Their opposed actions define an antagonistic pair.
\item \textbf{False.} Sustainable development means exploiting resources at a rate that allows them to regenerate and that does not compromise the needs of future generations — not stopping all exploitation, which would make economic and social development impossible.
\end{enumerate}
\end{corrige}

\begin{corrige}[Exercise 3 — Key definitions]
\begin{enumerate}
\item \textbf{Synovial joint:} a freely movable joint in which the bone ends, covered with articular cartilage, are separated by a fluid-filled synovial cavity and enclosed by a fibrous joint capsule lined with a synovial membrane; movement is stabilised by ligaments.
\item \textbf{Sarcomere:} the contractile unit of a myofibril, extending from one Z-line to the next, composed of overlapping thick (myosin) and thin (actin) filaments whose relative sliding produces contraction.
\item \textbf{Oxygen debt:} the extra volume of oxygen that must be taken in after strenuous exercise, above the resting requirement, to oxidise the accumulated lactic acid (to \ce{CO2} and water, or back to glucose in the liver) and to restore ATP and creatine phosphate reserves.
\item \textbf{Sustainable development:} development that meets the needs of the present generation without compromising the ability of future generations to meet their own needs, by balancing environmental protection, economic viability and social equity.
\end{enumerate}
\end{corrige}

\begin{corrige}[Exercise 4 — Direct calculations]
\begin{enumerate}
\item Decrease in length $= 2.2 - 1.6 = 0.6\ \mu\mathrm{m}$.
\[ \text{Percentage decrease} = \frac{0.6}{2.2}\times 100 = 27.3\% \]
\item Shortening per sarcomere $= 0.6\ \mu\mathrm{m}$. For $1\,500$ sarcomeres in series, the shortenings add:
\[ 1\,500 \times 0.6 = 900\ \mu\mathrm{m} = 0.9\ \mathrm{mm} \]
The fibre shortens by \textbf{0.9 mm}.
\item Increase $= 8.5 - 1.0 = 7.5\ \mathrm{mmol\,dm^{-3}}$.
\[ \text{Multiple of resting value} = \frac{8.5}{1.0} = 8.5 \]
The lactic acid concentration rose by $7.5\ \mathrm{mmol\,dm^{-3}}$, i.e. to \textbf{8.5 times} the resting value.
\end{enumerate}
\end{corrige}

\begin{corrige}[Exercise 5 — Labelling the human skeleton]
\begin{enumerate}
\item The leader lines should be labelled: \textbf{A} cranium; \textbf{B} vertebral column; \textbf{C} rib cage; \textbf{D} humerus; \textbf{E} radius and ulna; \textbf{F} pelvic girdle; \textbf{G} femur; \textbf{H} tibia and fibula.
\item One function of each region (any one valid function accepted):
\begin{itemize}
\item Cranium — \textbf{protection} of the brain;
\item Vertebral column — \textbf{support} of the trunk and protection of the spinal cord;
\item Rib cage — \textbf{protection} of the heart and lungs (and movement in breathing);
\item Humerus — \textbf{movement}: lever for the muscles of the upper arm;
\item Radius and ulna — \textbf{movement}: levers allowing rotation and flexion of the forearm;
\item Pelvic girdle — \textbf{support} of body weight and protection of abdominal organs;
\item Femur — \textbf{support} and \textbf{movement}: main weight-bearing lever of the thigh;
\item Tibia and fibula — \textbf{support} and \textbf{movement}: levers of the lower leg.
\end{itemize}
(The long bones also contain red bone marrow, the site of \textbf{blood cell production}.)
\item The joint between the femur and the pelvic girdle (the hip joint) is a \textbf{ball-and-socket joint}: the head of the femur fits into the acetabulum of the pelvis.
\end{enumerate}
\end{corrige}

\begin{corrige}[Exercise 6 — Comparing joints]
\begin{enumerate}
\item Comparison table:
\begin{center}
\begin{tabularx}{\linewidth}{|l|X|X|}
\hline
\rowcolor{popblueL} Feature & Hinge joint & Ball-and-socket joint \\
\hline
Example & Elbow (humerus--ulna), knee & Shoulder, hip \\
\hline
Planes of movement & One plane only & All planes (three) \\
\hline
Types of movement & Flexion and extension only & Flexion, extension, abduction, adduction, rotation, circumduction \\
\hline
\end{tabularx}
\end{center}
\item The glenoid cavity of the scapula is \textbf{shallow} and the head of the humerus is large relative to it, so the bony fit is loose; stability depends mainly on muscles and ligaments rather than on the shape of the socket. In contrast, the ulna fits deeply into the trochlea of the humerus, giving the elbow great bony stability. Hence the shoulder, which sacrifices stability for mobility, dislocates more easily.
\item The bones are held together by \textbf{ligaments} (tough, slightly elastic bands of fibrous connective tissue joining bone to bone) and by the \textbf{fibrous joint capsule}. Their role is to \textbf{stabilise the joint and prevent dislocation} while still permitting movement; the capsule also encloses the synovial cavity.
\end{enumerate}
\end{corrige}

\begin{corrige}[Exercise 7 — Smooth movement at a synovial joint]
\begin{enumerate}
\item The diagram must show: two bone ends covered by a layer of \textbf{articular cartilage}; a \textbf{synovial membrane} lining the inner face of the capsule; the \textbf{synovial cavity} containing \textbf{synovial fluid}; an outer fibrous \textbf{joint capsule}; and \textbf{ligaments} bridging the two bones outside the capsule. Labels should be placed with leader lines, a title given, and the drawing large and in pencil (GCE drawing conventions).
\begin{popfigure}
\begin{tikzpicture}[scale=1.2, every node/.style={font=\small}]
% Bone 1 (proximal)
\draw[fill=popblueL!30, draw=popdark, thick] (-3,0) rectangle (-1,2);
\draw[fill=white, draw=popblue, thick] (-3,2) -- (-2.8,2.2) -- (-1.2,2.2) -- (-1,2) -- cycle;
\draw[fill=poppink!50, draw=poppink, thick] (-3,2) rectangle (-1,2.3);
\node[align=center] at (-2,1) {Bone};
\node[align=center] at (-2,2.15) {Articular\\cartilage};

% Synovial cavity
\draw[fill=popblue!20, draw=popblue, thick] (-1,0.3) rectangle (1,1.7);
\node[align=center, text width=1.5cm] at (0,1) {Synovial cavity\\with synovial fluid};

% Bone 2 (distal)
\draw[fill=popblueL!30, draw=popdark, thick] (1,0) rectangle (3,2);
\draw[fill=white, draw=popblue, thick] (1,2) -- (1.2,2.2) -- (2.8,2.2) -- (3,2) -- cycle;
\draw[fill=poppink!50, draw=poppink, thick] (1,2) rectangle (3,2.3);
\node[align=center] at (2,1) {Bone};
\node[align=center] at (2,2.15) {Articular\\cartilage};

% Joint capsule (fibrous outer layer)
\draw[draw=popdark, thick, dashed] (-3,0) -- (-3,-0.3) -- (3,-0.3) -- (3,0);
\draw[draw=popdark, thick, dashed] (-3,2.3) -- (-3,2.6) -- (3,2.6) -- (3,2.3);
\draw[draw=popdark, thick, dashed] (-3,-0.3) -- (-3,2.6);
\draw[draw=popdark, thick, dashed] (3,-0.3) -- (3,2.6);
\node[align=center] at (-3.7,1.15) {Fibrous joint\\capsule};

% Synovial membrane (inner lining)
\draw[draw=popgreen, thick] (-2.9,0.1) -- (-2.9,2.4) -- (2.9,2.4) -- (2.9,0.1) -- cycle;
\node[align=center, text width=1.5cm] at (-3.7,0.5) {Synovial\\membrane};

% Ligaments
\draw[draw=poporange, ultra thick] (-2.5,-0.3) -- (-2.5,-0.6) -- (2.5,-0.6) -- (2.5,-0.3);
\draw[draw=poporange, ultra thick] (-2.5,2.6) -- (-2.5,2.9) -- (2.5,2.9) -- (2.5,2.6);
\node[align=center] at (0,-0.7) {Ligaments};
\node[align=center] at (0,3.0) {Ligaments};
\end{tikzpicture}
\legende{Structure of a typical synovial joint (e.g. knee or elbow).}
\end{popfigure}

\item Contribution of each structure:
\begin{itemize}
\item \textbf{Articular cartilage:} smooth, glassy covering of the bone ends that provides a low-friction gliding surface and absorbs shock;
\item \textbf{Synovial membrane:} secretes synovial fluid into the cavity;
\item \textbf{Synovial fluid:} lubricates the cartilage surfaces (reducing friction to almost zero) and nourishes the cartilage, which has no blood supply;
\item \textbf{Joint capsule:} encloses and seals the joint, retaining the synovial fluid;
\item \textbf{Ligaments:} hold the bones in correct alignment, preventing dislocation while allowing movement.
\end{itemize}
\item With too little synovial fluid, lubrication fails: the cartilage surfaces rub directly against each other, so friction rises sharply. This wears away the cartilage, irritates nerve endings (causing \textbf{pain}) and makes the joint surfaces resist gliding (causing \textbf{stiffness}), as occurs in some forms of arthritis.
\end{enumerate}
\end{corrige}

\begin{corrige}[Exercise 8 — Interpreting lactic acid data]
\begin{enumerate}
\item The graph should have \textbf{time (min) on the horizontal axis} and \textbf{lactic acid concentration ($\mathrm{mmol\,dm^{-3}}$) on the vertical axis}, with both axes labelled and scaled, points plotted at (0; 1.0), (2; 3.5), (4; 7.8), (6; 8.2), (10; 6.0), (20; 2.5), (30; 1.1), and a smooth curve drawn through them. The curve rises steeply to a peak near minute 6, then falls gradually back to the resting level by minute 30.
\begin{popfigure}
\begin{tikzpicture}
\begin{axis}[
    width=\linewidth,
    height=8cm,
    xlabel={Time (min)},
    ylabel={Lactic acid concentration (mmol dm$^{-3}$)},
    xmin=0, xmax=32,
    ymin=0, ymax=9,
    xtick={0,2,4,6,10,20,30},
    ytick={0,1,2,3,4,5,6,7,8},
    grid=major,
    grid style={popline},
    tick label style={font=\small},
    label style={font=\small},
    legend style={at={(0.98,0.98)}, anchor=north east, font=\small},
]
\addplot[popblue, thick, mark=*, mark size=2] coordinates {
    (0,1.0) (2,3.5) (4,7.8) (6,8.2) (10,6.0) (20,2.5) (30,1.1)
};
\addlegendentry{Blood lactic acid}
\end{axis}
\end{tikzpicture}
\legende{Blood lactic acid concentration during and after a 4-minute sprint.}
\end{popfigure}

\item Between minutes 0 and 4 the sprint demands ATP faster than aerobic respiration can supply it: oxygen cannot be delivered to the muscles quickly enough. The muscle fibres therefore respire \textbf{anaerobically}, converting glucose to \textbf{lactic acid}, which accumulates and diffuses into the blood, so blood lactic acid rises from 1.0 to $7.8\ \mathrm{mmol\,dm^{-3}}$.
\item After the sprint ends (minute 4), lactic acid already produced is still present in the muscles and blood; its removal is not instantaneous because it must be transported to the liver and oxidised, which itself requires oxygen and time. Hence the concentration peaks around minute 6 and only returns to near resting level after about 30 minutes.
\item The \textbf{oxygen debt} is the extra oxygen consumed after exercise, above the resting rate, needed to oxidise the accumulated lactic acid and to restore ATP and creatine phosphate stores. On the graph it corresponds to the whole recovery period (minutes 4--30) during which lactic acid is being cleared. The student continues to \textbf{breathe deeply} to take in this extra oxygen and to expel the \ce{CO2} produced by the oxidation of lactic acid.
\end{enumerate}
\end{corrige}

\begin{corrige}[Exercise 9 — Structured question: support and movement (20 marks)]
\textbf{Indicative mark scheme and examiner expectations:}
\begin{enumerate}
\item \textbf{(3 marks)} Any three of: \textbf{support} of the body / maintenance of shape; \textbf{protection} of delicate organs (brain, heart, lungs, spinal cord); \textbf{movement} — bones act as levers for muscle attachment; \textbf{blood cell production} in the red bone marrow; storage of minerals (calcium, phosphorus). \emph{1 mark each, maximum 3.}
\item \textbf{(4 marks)} In order from the pectoral girdle: \textbf{humerus} $\rightarrow$ \textbf{radius and ulna} $\rightarrow$ \textbf{carpals} $\rightarrow$ \textbf{metacarpals} $\rightarrow$ \textbf{phalanges}. \emph{1 mark per correct bone group in the correct order; the sequence must run proximally to distally.}
\item \textbf{(6 marks)} Structure: bone ends covered with \textbf{articular cartilage} (1); joint enclosed by a fibrous \textbf{capsule} lined with a \textbf{synovial membrane} (1); \textbf{synovial cavity} containing \textbf{synovial fluid} (1); bones held by \textbf{ligaments} (1). Friction minimised: smooth cartilage surfaces glide over each other (1); synovial fluid lubricates the surfaces, so friction is almost eliminated (1).
\item \textbf{(5 marks)} The \textbf{A-band} is the length of the thick myosin filaments, which do not change length (1). During contraction, the myosin heads pull the thin actin filaments \textbf{inwards towards the centre} of the sarcomere (sliding filament theory) (1). The Z-lines are therefore drawn closer together (1); the \textbf{I-band} (actin only) and the \textbf{H-zone} (myosin only) are progressively overlapped by the sliding filaments and so \textbf{shorten} (1), while the A-band stays constant because filament length is unchanged (1).
\item \textbf{(2 marks)} \textbf{Calcium ions (\ce{Ca^{2+}}):} released from the sarcoplasmic reticulum, they bind to troponin and expose the myosin-binding sites on actin, allowing cross-bridges to form (1). \textbf{ATP:} provides the energy for the power stroke and, on binding to the myosin head, causes its \textbf{detachment} from actin so the cycle can repeat (1).
\end{enumerate}
\emph{Examiner expectations: precise terminology (sarcomere, Z-line, cross-bridge), correct bone order, and explanations — not mere lists — for parts (c) and (d).}
\end{corrige}

\begin{corrige}[Exercise 10 — Essay: muscles, antagonism and levers (20 marks)]
\textbf{Model plan and indicative mark scheme:}

\textbf{Introduction (2 marks).} A skeletal muscle exerts force only by shortening: contraction can pull a bone towards a joint, but a muscle cannot actively lengthen to push a bone away. The statement is therefore correct, and the body solves this limitation through antagonistic pairs and levers.

\textbf{Mechanism of contraction (6 marks).} Describe the sliding filament theory: nerve impulse $\rightarrow$ release of \ce{Ca^{2+}} from the sarcoplasmic reticulum $\rightarrow$ exposure of binding sites on actin $\rightarrow$ myosin heads attach, pivot (power stroke, ATP hydrolysed) and detach $\rightarrow$ actin filaments slide inwards, Z-lines approach, the sarcomere \emph{shortens}. Force is thus generated exclusively by \textbf{tension} (pulling); relaxation is passive and lengthening is produced by an external force, never by the muscle itself. \emph{Expect: cross-bridge cycle, roles of \ce{Ca^{2+}} and ATP, shortening of I-band/H-zone.}

\textbf{Antagonistic pairs — the human arm (6 marks).} Flexion of the elbow: the \textbf{biceps} contracts and pulls the radius upwards while the \textbf{triceps} relaxes and is passively stretched. Extension: the \textbf{triceps} contracts and pulls the ulna (olecranon) backwards while the biceps relaxes. Each movement is produced by a pull; the return movement is produced by the pull of the opponent muscle. Other examples accepted (quadriceps/hamstrings at the knee). \emph{Expect: correct identification of prime mover and antagonist in each movement.}

\textbf{Bones as levers (4 marks).} Bones act as levers rotating about joints (fulcra), with muscles supplying the effort through tendons. At the elbow (a third-order lever): fulcrum = elbow joint; effort = biceps insertion on the radius; load = weight in the hand. This arrangement gives a \textbf{speed/range advantage}: a small shortening of the muscle produces a large, rapid movement of the hand.
\begin{popfigure}
\begin{tikzpicture}[scale=1.1, every node/.style={font=\small}]
% Humerus
\draw[fill=popblueL!30, draw=popdark, thick] (0,0) -- (0,3) -- (0.5,3) -- (0.5,0) -- cycle;
\node at (0.25,1.5) {Humerus};
% Radius/Ulna (forearm)
\draw[fill=popblueL!30, draw=popdark, thick] (0,0) -- (-1.5,-2.5) -- (-1,-2.5) -- (0.5,0) -- cycle;
\node at (-0.75,-1.25) {Radius/Ulna};
% Fulcrum (elbow joint)
\draw[fill=poporange, draw=popdark] (0,0) circle (0.15);
\node[below right] at (0,0) {Fulcrum};
% Biceps (effort)
\draw[popgreen, thick, ->] (0.5,2.5) -- (0.5,1);
\node[right] at (0.5,1.75) {Effort (biceps)};
% Load (weight in hand)
\draw[popred, thick, ->] (-1.25,-2.5) -- (-1.25,-3.2);
\node[below] at (-1.25,-3.2) {Load};
% Triceps (antagonist, relaxed)
\draw[popmuted, thick, dashed] (-0.5,2.5) -- (-0.5,0.5);
\node[left] at (-0.5,1.5) {Triceps (relaxed)};
\end{tikzpicture}
\legende{Elbow joint as a third-order lever during flexion.}
\end{popfigure}

\textbf{Diagram and conclusion (2 marks).} A clearly labelled diagram of the arm showing biceps, triceps, elbow joint (fulcrum), effort and load. Conclusion: because muscles only pull, coordinated movement requires at least two muscles working antagonistically across a joint acting as a lever system.

\emph{Examiner expectations: logical essay structure, accurate terminology, at least one labelled diagram, and explicit linkage of the three required themes to the statement.}
\end{corrige}

\begin{corrige}[Exercise 11 — Essay: sustainable development in Cameroon (20 marks)]
\textbf{Indicative mark scheme and examiner expectations:}
\begin{enumerate}
\item \textbf{(4 marks)} Definition (2): development that meets the needs of the present \textbf{without compromising the ability of future generations to meet their own needs}. Three pillars (2): \textbf{environmental} (protection of ecosystems and resources), \textbf{economic} (viable, durable economic activity) and \textbf{social} (equity, health, education, well-being of communities).
\item \textbf{(8 marks)} Example using the \textbf{forests of the South Region} (any coherent named resource accepted):
\begin{itemize}
\item Present needs: timber provides income, employment and export earnings; forest land supports farming and fuelwood collection (2);
\item Environmental measures: selective logging rather than clear-felling; replanting/afforestation; respecting minimum cutting diameters so trees can regenerate; protecting watersheds and biodiversity (2);
\item Economic measures: local processing of timber to add value; certification and control of illegal logging so revenue is durable; developing ecotourism and non-timber forest products (2);
\item Social measures: involving local communities in management; guaranteeing their access to fuelwood and food; using revenues for schools and health centres, so that future generations inherit both the resource and the means to use it (2).
\end{itemize}
\item \textbf{(8 marks)} Interdependence in a rural community, with concrete examples:
\begin{itemize}
\item Environment $\rightarrow$ economy: fertile soils and regular rainfall (Western Highlands) allow coffee and food-crop production; soil erosion from over-cultivation destroys yields and income (2);
\item Economy $\rightarrow$ social: income from cocoa or fishing pays school fees and health care; poverty, conversely, forces over-exploitation (cutting firewood, overfishing) (2);
\item Social $\rightarrow$ environment: education and community organisation enable sustainable practices (tree planting, closed fishing seasons); population pressure without alternatives degrades land (2);
\item Conclusion: the three pillars form a cycle — weakening one undermines the others; true sustainability requires all three to be pursued together (2).
\end{itemize}
\end{enumerate}
\emph{Examiner expectations: an accurate definition, one clearly named Cameroonian resource treated in depth, concrete and realistic local examples, and an explicit demonstration that the three pillars depend on one another.}
\end{corrige}

\newpage

\coursec{Summary sheet}

\begin{popfigure}
\begin{tikzpicture}[mindmap, grow cyclic, every node/.style={concept, font=\small, text=white},
  root concept/.append style={concept color=popdark, font=\small\bfseries, minimum size=2.6cm},
  level 1/.append style={level distance=3.4cm, sibling angle=60},
  level 1 concept/.append style={minimum size=2.2cm}]
\node[root concept]{Support \& Locomotion}
  child[concept color=popblue]{ node {Types of skeletons}
    child[concept color=popblue]{ node[concept, minimum size=1.7cm, font=\scriptsize, text=white]{Hydrostatic, exo-, endoskeleton} }}
  child[concept color=popgreen]{ node {Vertebrate skeleton}
    child[concept color=popgreen]{ node[concept, minimum size=1.7cm, font=\scriptsize, text=white]{Axial \& appendicular; bone structure} }}
  child[concept color=poporange]{ node {Joints}
    child[concept color=poporange]{ node[concept, minimum size=1.7cm, font=\scriptsize, text=white]{Fibrous, cartilaginous, synovial} }}
  child[concept color=poppink]{ node {Smooth movement}
    child[concept color=poppink]{ node[concept, minimum size=1.7cm, font=\scriptsize, text=white]{Synovial fluid, cartilage, ligaments} }}
  child[concept color=popgold]{ node {Muscle contraction}
    child[concept color=popgold]{ node[concept, minimum size=1.7cm, font=\scriptsize, text=white]{Sliding filament theory; ATP} }}
  child[concept color=popmuted]{ node {Sustainable development}
    child[concept color=popmuted]{ node[concept, minimum size=1.7cm, font=\scriptsize, text=white]{Needs of today without compromising tomorrow} }};
\end{tikzpicture}
\legende{Mind map of Chapter BIO05: Support and locomotion, with the link to sustainable development.}
\end{popfigure}

\begin{bilan}
\textbf{1. Key definitions}
\begin{itemize}
\item \cle{Skeleton}: rigid or fluid framework that supports the body, protects organs and provides attachment for muscles.
\item \cle{Hydrostatic skeleton}: fluid-filled cavity surrounded by muscles (earthworm); \cle{exoskeleton}: external cuticle of chitin (arthropods); \cle{endoskeleton}: internal skeleton of bone and cartilage (vertebrates).
\item \cle{Bone}: living tissue made of collagen fibres hardened by calcium phosphate and calcium carbonate salts; contains osteocytes, blood vessels and marrow.
\item \cle{Axial skeleton}: skull, vertebral column, ribs and sternum. \cle{Appendicular skeleton}: limb bones plus pectoral and pelvic girdles.
\item \cle{Joint (articulation)}: point where two bones meet. \cle{Ligament} joins bone to bone; \cle{tendon} joins muscle to bone.
\item \cle{Antagonistic muscles}: pairs producing opposite movements (e.g. biceps flexes, triceps extends the forearm).
\item \cle{Sarcomere}: contractile unit of a myofibril, between two Z-lines.
\item \cle{Sustainable development}: development that meets present needs without compromising the ability of future generations to meet their own needs.
\end{itemize}

\textbf{2. Facts and structures to know}
\begin{itemize}
\item Vertebral column regions: cervical (7), thoracic (12), lumbar (5), sacral (5 fused), caudal/coccygeal (4 fused) vertebrae. The \cle{atlas} supports the skull; the \cle{axis} allows rotation of the head.
\item A typical vertebra: centrum (body), neural arch and spine (protecting the spinal cord), transverse processes, facets for articulation.
\item Synovial joint components: articular \cle{cartilage} (reduces friction), \cle{synovial membrane} secreting \cle{synovial fluid} (lubrication and nourishment), joint capsule, ligaments (stability).
\item Joint types: fibrous (immovable, skull sutures), cartilaginous (slightly movable, between vertebrae), synovial (freely movable: ball-and-socket at hip/shoulder; hinge at knee/elbow; pivot at atlas--axis; gliding in the wrist).
\item Muscle hierarchy: muscle $\rightarrow$ fascicles $\rightarrow$ fibres (cells) $\rightarrow$ myofibrils $\rightarrow$ myofilaments: thick \cle{myosin} and thin \cle{actin} (with tropomyosin and troponin).
\item Sarcomere bands: A band (myosin zone), I band (actin only), H zone; during contraction the I band and H zone \textbf{shorten}, the A band stays \textbf{constant}.
\end{itemize}

\textbf{3. The sliding filament theory (method to explain contraction)}
\begin{enumerate}
\item A nerve impulse arrives at the neuromuscular junction; \ce{Ca^{2+}} is released from the sarcoplasmic reticulum.
\item \ce{Ca^{2+}} binds troponin, moving tropomyosin and exposing myosin-binding sites on actin.
\item Myosin heads (charged with ADP $+$ P$_i$) attach to actin, forming \cle{cross-bridges}.
\item Power stroke: heads tilt, pulling actin towards the centre of the sarcomere (ADP and P$_i$ released).
\item ATP binds myosin, breaking the cross-bridge; ATP is hydrolysed, re-cocking the head. The cycle repeats while \ce{Ca^{2+}} and ATP are present.
\item When stimulation stops, \ce{Ca^{2+}} is pumped back, binding sites are covered, and the muscle relaxes.
\end{enumerate}
Roles of ATP: detachment of myosin heads, re-cocking, and pumping \ce{Ca^{2+}} back. Rigor mortis occurs because without ATP, cross-bridges cannot detach.

\textbf{4. Sustainable development essentials}
\begin{itemize}
\item Three pillars: \cle{economic}, \cle{social} and \cle{environmental}.
\item Principles: rational use of resources, recycling, renewable energy, conservation of biodiversity (e.g. Korup National Park, Mount Cameroon forests), reforestation.
\item Link to the chapter: overexploitation of forests and bushmeat threatens species and ecosystems; sustainable practices protect livelihoods for future generations.
\end{itemize}

\textbf{5. Common mistakes}
\begin{itemize}
\item Confusing \textbf{tendon} (muscle--bone) with \textbf{ligament} (bone--bone).
\item Saying the A band shortens: only the I band and H zone shorten.
\item Claiming ATP is used for the power stroke: ATP is needed to \textbf{detach} and re-cock myosin heads.
\item Writing that muscles ``push'': muscles can only \textbf{pull} (contract); hence antagonistic pairs.
\item Mixing up the atlas (supports skull, ``yes'' movement) and the axis (rotation, ``no'' movement).
\end{itemize}

\textbf{6. GCE tips}
\begin{itemize}
\item In diagrams, always label: cartilage, synovial fluid, capsule, ligament, tendon, and the antagonistic muscles.
\item For essay questions on contraction, present the sliding filament theory as ordered steps and mention \ce{Ca^{2+}}, troponin, cross-bridges and ATP explicitly.
\item Quote precise numbers of vertebrae; examiners reward accuracy.
\item For sustainable development questions, define the concept, state the three pillars, and give a concrete Cameroonian example.
\end{itemize}
\end{bilan}

\findecours

\end{document}
